\chapter{THE ONSET OF THERMAL INSTABILITY IN FLUID SPHERES AND SPHERICAL SHELLS}
\label{ch:6}

\section{Introduction}\label{sec:55}
In this chapter we shall extend the considerations of the preceding
chapters to a spherical geometry and to problems relating to the thermal
instability of fluid spheres and spherical shells with internal heat sources.
Apart from the general interest, from the point of view of applied
mathematics, which attaches to such extensions of a basic physical
theory, as the nature and origin of thermal instability, the extensions
are also of interest to geophysics in several connexions: for example,
with the pattern of the convective motions which might prevail in the
fluid core of the earth and the relation of such motions to theories con\-cerning the origin of the earth's magnetic field and its secular variations.
There are other geophysical connexions in which the nature of convection
in a spherical geometry might also be important: in connexion with the
question whether the roughly equal divisions of the earth's surface into a
land and an ocean hemisphere might be related to a convection pattern
in an original fluid earth; and in connexion also with the question whether
a slow, continuing, large-scale convection in the earth's mantle might be
related to the origin and the movements of the continents. The relevance
of these latter questions in the particular geophysical and geological
contexts are matters of controversy at the present time; but they do not
affect the interest in the problems \emph{per se}.

\section{The perturbation equations}\label{sec:56}
Consider a spherical shell of incompressible fluid subject to a spheri\-cally symmetric radial gravitational field $g(\boldsymbol{r})$, where $g(\boldsymbol{r})$ is a function
of $r$ only. In any specific case, $g(\boldsymbol{r})$ will be a known function of $r$. Thus,
if we should consider a homogeneous fluid sphere of density $\rho$,
\begin{equation}\label{eq:6-1}
g(r) = \tfrac{4}{3}\pi\rho G r = \gamma r,\quad\text{say},
\end{equation}
where $G$ is the constant of gravitation. Similarly, if we should consider
a spherical shell of density $\rho$ overlying a core of radius $R_1$ and mass $M_0$,
\begin{equation}\label{eq:6-2}
g(r) = \tfrac{G}{r^2}\bigl(M_0+\tfrac{4}{3}\pi\rho[r^3-R_1^3]\bigr).
\end{equation}

We shall suppose that there is a distribution of heat sources $\epsilon$ which
maintains a radial temperature gradient in the fluid. The gradient will
be determined by the equation of heat conduction,
\begin{equation}\label{eq:6-3}
\kappa\nabla^2 T = -\epsilon.
\end{equation}
If $\epsilon$ and $\kappa$ are assumed to be constants, equation~\eqref{eq:6-3} leads to the tempera\-ture distribution:
\begin{equation}\label{eq:6-4}
T = \bar{T}_0 - \bar{\beta}_0 r^2 + \bar{\beta}_2,
\end{equation}
where
\begin{equation}\label{eq:6-5}
\bar{\beta}_0 = \epsilon/6\kappa,
\end{equation}
and $\bar{T}_0$ and $\bar{\beta}_2$ are constants. By considering the stability of this initial
state, we readily find that the perturbation equations in the Boussinesq
approximation are
\begin{equation}\label{eq:6-6}
\frac{\partial u_i}{\partial t} = -\frac{1}{\rho}\frac{\partial \delta p}{\partial x_i} + g\alpha\Theta\frac{x_i}{r} + \nu\nabla^2 u_i,
\end{equation}
and
\begin{equation}\label{eq:6-7}
\frac{\partial \Theta}{\partial t} = \beta(u_j x_j/r) + \kappa\nabla^2\Theta,
\end{equation}
where $\alpha$ is the coefficient of volume expansion, and $\delta p$ and $\Theta$ are the
perturbations in the pressure and the temperature, respectively.
According to equation~\eqref{eq:6-4}, we may write
\begin{equation}\label{eq:6-8}
\beta = 2\bar{\beta}_0,
\end{equation}
and
\begin{equation}\label{eq:6-9}
\gamma = g/r.
\end{equation}
Equations~\eqref{eq:6-6} and \eqref{eq:6-7} may now be rewritten in the forms:
\begin{equation}\label{eq:6-10}
\frac{\partial u_i}{\partial t} = -\frac{1}{\rho}\frac{\partial \delta p}{\partial x_i} + \psi\Theta x_i + \nu\nabla^2 u_i,
\end{equation}
and
\begin{equation}\label{eq:6-11}
\frac{\partial \Theta}{\partial t} = \beta(u_j x_j/r) + \kappa\nabla^2\Theta,
\end{equation}
where
\begin{equation}\label{eq:6-12}
\psi(r) = \alpha\gamma(r).
\end{equation}
In addition to equations~\eqref{eq:6-10} and \eqref{eq:6-11}, we have the equation of con\-tinuity
\begin{equation}\label{eq:6-13}
\frac{\partial u_i}{\partial x_i} = 0.
\end{equation}

Following our standard practice, we eliminate the term $\operatorname{grad}(\delta p/\rho)$
in equation~\eqref{eq:6-10} by taking its curl: We obtain
\begin{equation}\label{eq:6-14}
\frac{\partial \omega_i}{\partial t} = \psi\Bigl(r\frac{\partial \Theta}{\partial x_i} - \frac{x_i}{r}\frac{\partial(r\Theta)}{\partial r}\Bigr) + \Theta x_i\frac{1}{r}\frac{d\psi}{dr} + \nu\nabla^2\omega_i,
\end{equation}
where $\omega_i$ is the vorticity. Taking the curl of this equation once again,
we obtain
\begin{equation}\label{eq:6-15}
\frac{\partial}{\partial t}\nabla^2 u_i = -\mathscr{Q}_i\,\Theta + \psi\nabla^2(\Theta x_i) + \nu\nabla^4 u_i,
\end{equation}
where $\mathscr{Q}_i$ stands for the differential operator
\begin{equation}\label{eq:6-16}
\mathscr{Q}_i = \frac{\partial}{\partial x_i}\Bigl(\psi r\frac{\partial}{\partial r} + \psi + r\frac{d\psi}{dr}\Bigr) - \psi\frac{\partial}{\partial x_i}.
\end{equation}
On expanding, we find
\begin{equation}\label{eq:6-17}
\mathscr{Q}_i = \frac{\partial}{\partial x_i}\Bigl(\psi r\frac{\partial}{\partial r} + r\frac{d\psi}{dr}\Bigr).
\end{equation}
Now it can be directly verified that
\begin{equation}\label{eq:6-18}
\mathscr{Q}_i\nabla^2 u_i = \nabla^2(\mathscr{Q}_i u_i),\quad\text{if } \boldsymbol{u} \text{ is solenoidal;}
\end{equation}
and since $\boldsymbol{u}$ and $\boldsymbol{\omega}$ are both solenoidal, we obtain from equations
\eqref{eq:6-14} and \eqref{eq:6-15}:
\begin{equation}\label{eq:6-19}
\frac{\partial \omega_i}{\partial t} = \psi\Bigl(r\frac{\partial \Theta}{\partial x_i} - x_i\frac{\partial \Theta}{\partial r}\Bigr) + \nu\nabla^2\omega_i,
\end{equation}
and
\begin{equation}\label{eq:6-20}
\frac{\partial}{\partial t}\nabla^2 u_i = \nabla^2(\psi\Theta x_i) - \frac{\partial}{\partial x_i}\Bigl(\psi r\frac{\partial \Theta}{\partial r} + r\frac{d\psi}{dr}\Theta\Bigr) + \nu\nabla^4 u_i,
\end{equation}
where
\begin{equation}\label{eq:6-21}
\gamma(r) = \psi r\frac{\partial}{\partial r}\Theta + r\frac{d\psi}{dr}\Theta.
\end{equation}

We can eliminate $\Theta$ from equation~\eqref{eq:6-20} by observing that the operators
$\nabla^2$ and $\mathscr{L}^2$ commute (see equation~\eqref{eq:6-25} below):

Therefore,
\begin{equation}\label{eq:6-22}
\frac{\partial}{\partial t}\bigl(\nabla^2 - \tfrac{1}{\nu}\tfrac{\partial}{\partial t}\bigr)(u_i x_i) = \psi\,\mathscr{L}^2\,\Theta + \nu\nabla^4(u_i x_i).
\end{equation}

Now making use of equation~\eqref{eq:6-11}, we obtain
\begin{equation}\label{eq:6-23}
\Bigl(\nabla^2-\frac{1}{\kappa}\frac{\partial}{\partial t}\Bigr)\Bigl(\nabla^2 - \frac{1}{\nu}\frac{\partial}{\partial t}\Bigr)\nabla^2(u_i x_i) = -\frac{\psi}{\kappa\nu}\,\mathscr{L}^2\,\beta(u_j x_j/r).
\end{equation}

Equations~\eqref{eq:6-11}, \eqref{eq:6-19}, \eqref{eq:6-20}, and \eqref{eq:6-23} are the basic equations of the
problem.

\subsection{The operator $\mathscr{L}^2$}\label{sec:56a}
For future reference we may note here certain elementary properties
of the operator $\mathscr{L}^2$. In spherical polar coordinates $r$, $\theta$, and $\varphi$,
\begin{equation}\label{eq:6-25}
\mathscr{L}^2 = -\frac{1}{\sin\theta}\frac{\partial}{\partial\theta}\Bigl(\sin\theta\frac{\partial}{\partial\theta}\Bigr) - \frac{1}{\sin^2\theta}\frac{\partial^2}{\partial\varphi^2};
\end{equation}
$\mathscr{L}^2$ is, indeed, the operator representing the square of the angular
momentum. The characteristic values are, therefore, $l(l+1)$; and the
eigenfunctions belonging to them are the spherical harmonics $Y_l^m(\theta,\varphi)$:
\begin{equation}\label{eq:6-26}
\mathscr{L}^2 Y_l^m(\theta,\varphi) = l(l+1)Y_l^m(\theta,\varphi),
\end{equation}
where
\begin{equation}\label{eq:6-27}
Y_l^m(\theta,\varphi) = P_l^m(\cos\theta)e^{im\varphi},
\end{equation}
and $P_l^m(\cos\theta)$ are the associated Legendre polynomials. The corre\-sponding normalization integral is
\begin{equation}\label{eq:6-28}
\int_0^{2\pi}\!\!\int_0^{\pi} |Y_l^m(\theta,\varphi)|^2 \sin\theta\,d\theta\,d\varphi = N_l^m,
\end{equation}
where
\begin{equation}\label{eq:6-29}
N_l^m = \frac{4\pi}{2l+1}\frac{(l+m)!}{(l-m)!}.
\end{equation}

\subsection{The analysis into normal modes}\label{sec:56b}
Returning to equations~\eqref{eq:6-11}, \eqref{eq:6-19}, \eqref{eq:6-20}, and \eqref{eq:6-23}, we must, in accordance
with our general procedure, analyse the disturbance into normal modes.
In the present instance, the analysis in terms of spherical harmonics is
clearly indicated. Accordingly, we write
\begin{equation}\label{eq:6-30}
u_i x_i = u_r r = W(r) Y_l^m e^{pt},
\end{equation}
and
\begin{equation}\label{eq:6-31}
\Theta = G(r) Y_l^m e^{pt},
\end{equation}
Since
\begin{equation}\label{eq:6-32}
\mathscr{L}^2 Y_l^m = l(l+1) Y_l^m,
\end{equation}
the effect of $\nabla^2$ on a function which is a product of a function of $r$ only,
and a spherical harmonic $Y_l^m$ is given by
\begin{equation}\label{eq:6-33}
\nabla^2[f(r)Y_l^m] = (\mathscr{D}_l^2 f)\,Y_l^m,
\end{equation}
where
\begin{equation}\label{eq:6-34}
\mathscr{D}_l^2 = \frac{d^2}{dr^2} + \frac{2}{r}\frac{d}{dr} - \frac{l(l+1)}{r^2}.
\end{equation}
In view of this identity, the insertion of the forms of the solution given
in~\eqref{eq:6-30}, in equations~\eqref{eq:6-11}, \eqref{eq:6-19}, and \eqref{eq:6-23} gives
\begin{equation}\label{eq:6-35}
(\mathscr{D}_l^2 - p/\kappa)\,G = -\frac{\beta}{r}\,W,
\end{equation}
and
\begin{equation}\label{eq:6-36}
(\mathscr{D}_l^2 - p/\nu)\mathscr{D}_l^2\,W = \frac{\psi}{\nu}\,l(l+1)\,G,
\end{equation}
where we have measured $r$ in units of a suitable radius $R_1$ and let
\begin{equation}\label{eq:6-37}
\sigma = pR_1^2/\nu \quad\text{and}\quad p = \nu/R_1^2.
\end{equation}

\subsection{Boundary conditions}\label{sec:56c}
Solutions of equations~\eqref{eq:6-35}--\eqref{eq:6-36} must be sought which satisfy certain
boundary conditions. These are described below.

Let the spherical shell be confined between $r = 1$ and $\eta$. (The unit of
length chosen is then the outer radius $R_0$ of the shell.)
In all cases we must require that
\begin{equation}\label{eq:6-38}
W = 0\quad\text{and}\quad G = 0\quad\text{for}\quad r = 1 \quad\text{and}\quad r = \eta.
\end{equation}

The remaining boundary conditions on $W$ depend on the nature of the
surfaces at $r = 1$ and $\eta$. We shall consider two cases: the case when the
surface is rigid (as at the interface with the mantle of the liquid core of
the earth); and the case when the surface is free (as at the boundary of
an isolated fluid sphere in free space). On a rigid boundary, we must
require that the transverse components of the velocity
\begin{equation}\label{eq:6-39}
u_\theta \quad\text{and}\quad u_\varphi \quad\text{also vanish};
\end{equation}
while on a free boundary, we must require that the tangential viscous
stresses
\begin{equation}\label{eq:6-40}
p_{r\theta} \quad\text{and}\quad p_{r\varphi} \quad\text{vanish}.
\end{equation}

Consider first the implications of imposing the condition~\eqref{eq:6-39} on a
spherical bounding surface. From the equation of continuity,
\begin{equation}\label{eq:6-41}
\frac{1}{r^2}\frac{\partial}{\partial r}(r^2 u_r) + \frac{1}{r\sin\theta}\frac{\partial}{\partial\theta}(\sin\theta\, u_\theta) + \frac{1}{r\sin\theta}\frac{\partial u_\varphi}{\partial\varphi} = 0,
\end{equation}
it follows from the vanishing of $u_r$, $u_\theta$, and $u_\varphi$ on a surface, $r = \text{constant}$,
for all values of $\theta$ and $\varphi$, that
\begin{equation}\label{eq:6-42}
\frac{\partial u_r}{\partial r} = 0 \quad\text{(on a rigid spherical boundary)}.
\end{equation}
An equivalent form of this condition is
\begin{equation}\label{eq:6-43}
\frac{dW}{dr} = 0 \quad\text{(on a rigid spherical boundary)}.
\end{equation}

Consider next the implications of imposing the condition~\eqref{eq:6-40} on a
spherical bounding surface. The expressions for the viscous stresses
are
\begin{equation}\label{eq:6-44}
p_{r\theta} = \rho\nu\Bigl(r\frac{\partial}{\partial r}\frac{u_\theta}{r} + \frac{1}{r}\frac{\partial u_r}{\partial\theta}\Bigr),
\end{equation}
and
\begin{equation}\label{eq:6-45}
p_{r\varphi} = \rho\nu\Bigl(\frac{1}{r\sin\theta}\frac{\partial u_r}{\partial\varphi} + r\frac{\partial}{\partial r}\frac{u_\varphi}{r}\Bigr).
\end{equation}
It follows from the vanishing of $p_{r\theta}$ and $p_{r\varphi}$ on a surface, $r = \text{constant}$
(on which $u_r$ vanishes identically), that
\begin{equation}\label{eq:6-46}
r\frac{\partial}{\partial r}\Bigl(\frac{u_\theta}{r}\Bigr) = 0 \quad\text{and}\quad r\frac{\partial}{\partial r}\Bigl(\frac{u_\varphi}{r}\Bigr) = 0,
\end{equation}
and since $u_\theta$ vanishes on this surface, the condition~\eqref{eq:6-46} is equivalent to
\begin{equation}\label{eq:6-47}
\frac{\partial^2 W}{\partial r^2} = 0 \quad\text{(on a free spherical boundary)}.
\end{equation}
If we should be considering a complete fluid sphere, certain boundary
conditions must be satisfied at the centre which ensure that none of the
quantities have singularities there. From the differential equations
satisfied by $W$ and $G$, it follows that
\begin{equation}\label{eq:6-48}
W \quad\text{and}\quad G\quad\text{must behave like}\quad r^l \quad\text{as}\quad r \to 0.
\end{equation}

\subsection{The velocity field}\label{sec:56d}
After solving equations~\eqref{eq:6-35}--\eqref{eq:6-36} we must complete the solution by
specifying the entire velocity field; this requires the determination of
the remaining components of the velocity $u_\theta$ and $u_\varphi$ in terms of the radial
components of the velocity and the vorticity. The solution can best
be accomplished by making use of the following general representation
of an arbitrary solenoidal vector field. (A more detailed account
of this representation will be found in Appendix~III.)

Any solenoidal vector field can be expressed as a linear combination
of certain basic \emph{toroidal} (\textbf{T}) and \emph{poloidal} (\textbf{S}) fields given by
\begin{equation}\label{eq:6-49}
\mathbf{S} = \operatorname{curl}\operatorname{curl}(r\hat{\boldsymbol{r}}\,\mathscr{S}) \equiv -\frac{1}{r}\nabla^2_r(\mathscr{L}^2\mathscr{S})\hat{\boldsymbol{r}} - \frac{1}{r}\frac{\partial^2(r\mathscr{S})}{\partial r\,\partial\theta}\hat{\boldsymbol{\theta}} - \frac{1}{r\sin\theta}\frac{\partial^2(r\mathscr{S})}{\partial r\,\partial\varphi}\hat{\boldsymbol{\varphi}},
\end{equation}
and
\begin{equation}\label{eq:6-50}
\mathbf{T} = \operatorname{curl}(r\hat{\boldsymbol{r}}\,\mathscr{T}) \equiv -\frac{1}{\sin\theta}\frac{\partial \mathscr{T}}{\partial\varphi}\hat{\boldsymbol{\theta}} + \frac{\partial \mathscr{T}}{\partial\theta}\hat{\boldsymbol{\varphi}}.
\end{equation}

The essential properties of the vector fields $\mathbf{S}$ and $\mathbf{T}$, which are proved
in Appendix~III, may be summarized here:

\begin{enumerate}
\item[(i)] The fields $\mathbf{S}$ and $\mathbf{T}$ are both solenoidal.
\item[(ii)] For integration over a surface of a sphere of radius $r$ the fields
$\mathbf{S}$ and $\mathbf{T}$ have the following orthogonality properties.

If $\mathbf{S}$ and $\mathbf{S}'$ are derived from different spherical harmonics,
\begin{equation}\label{eq:6-51}
\int \mathbf{S}\cdot\mathbf{S}'\,d\Omega = 0,
\end{equation}
Similarly, if $\mathbf{T}$ and $\mathbf{T}'$ are derived from
different spherical harmonics,
\begin{equation}\label{eq:6-52}
\int \mathbf{T}\cdot\mathbf{T}'\,d\Omega = 0.
\end{equation}

\item[(iii)] If $\mathbf{S}$ and $\mathbf{S}'$ are derived from the same spherical harmonic,
\begin{equation}\label{eq:6-53}
\int \mathbf{S}\cdot\mathbf{S}'\,d\Omega = N_l^m l(l+1)\Bigl[\frac{1}{r^2}\Bigl(\frac{d(r\mathscr{S})}{dr}\Bigr)\Bigl(\frac{d(r\mathscr{S}')}{dr}\Bigr) + \frac{l(l+1)}{r^2}\mathscr{S}\,\mathscr{S}'\Bigr],
\end{equation}
where $N_l^m$ has the same meaning as in equation~\eqref{eq:6-29}. Similarly, if $\mathbf{T}$
and $\mathbf{T}'$ are derived from the same spherical harmonic,
\begin{equation}\label{eq:6-54}
\int \mathbf{T}\cdot\mathbf{T}'\,d\Omega = l(l+1)N_l^m\,\mathscr{T}\,\mathscr{T}'.
\end{equation}

\item[(iv)] Every poloidal field is orthogonal to every toroidal field:
\begin{equation}\label{eq:6-55}
\int \mathbf{S}\cdot\mathbf{T}\,d\Omega = 0.
\end{equation}

\item[(v)] curl $\mathbf{S}$ represents a toroidal field with the defining scalar
\begin{equation}\label{eq:6-56}
S = -\frac{l(l+1)}{r}\mathscr{S};
\end{equation}
and curl $\mathbf{T}$ represents a poloidal field with the defining scalar $\mathscr{T}$.

\item[(vi)] curl$^2\mathbf{S}$ is again a poloidal field with $S$ as its defining scalar;
similarly, curl$^2\mathbf{T}$ is a toroidal field with the defining scalar
\begin{equation}\label{eq:6-57}
S = -\frac{l(l+1)}{r}\frac{d\mathscr{T}}{dr}.
\end{equation}
\end{enumerate}

From the foregoing it follows that $W$ and $Z$, as we have defined them
in equation~\eqref{eq:6-30}, are related to the scalars $\mathscr{S}$ and $\mathscr{T}$ by
\begin{equation}\label{eq:6-58}
S = -\frac{l(l+1)}{r}\mathscr{S} \quad\text{and}\quad T = -\frac{l(l+1)}{r}\mathscr{T}.
\end{equation}

\section{The validity of the principle of the exchange of stabilities
for the case $\beta = \text{constant}$ and $\gamma = \text{constant}$}\label{sec:57}
In this section the discussion will be restricted to the case when both
$\beta$ and $\gamma$ are constants. Then (see equations~\eqref{eq:6-8}, \eqref{eq:6-9}, and \eqref{eq:6-12})
\begin{equation}\label{eq:6-59}
\beta = \beta_0 = \epsilon/6\kappa \quad\text{and}\quad \gamma = \gamma_0 = \tfrac{4}{3}\pi\rho G.
\end{equation}

Equations~\eqref{eq:6-35} and \eqref{eq:6-36} can then be combined to give
\begin{equation}\label{eq:6-60}
\mathscr{D}_l^2(\mathscr{D}_l^2 - \sigma)(\mathscr{D}_l^2 - p\sigma)W = l(l+1)\mathcal{C} W,
\end{equation}
where
\begin{equation}\label{eq:6-61}
\mathcal{C} = \frac{\alpha\beta\gamma R_1^4}{\kappa\nu}.
\end{equation}

We shall now show that when $\beta$ and $\gamma$ are constants, the principle
of the exchange of stabilities is valid. But first we shall prove certain
properties of the operator $\mathscr{D}_l^2$ which we shall find useful in the subsequent
reductions.

Consider the integral
\begin{equation}\label{eq:6-62}
\int_a^b \phi\,\frac{d}{dr}\Bigl[\frac{d}{dr}\Bigl(\frac{\psi}{r}\Bigr) - \frac{l(l+1)}{r^2}\psi\Bigr]r^2\,dr =
\int_a^b \phi\Bigl[\frac{d^2\psi}{dr^2} + \frac{2}{r}\frac{d\psi}{dr} - \frac{l(l+1)}{r^2}\psi\Bigr]r^2\,dr,
\end{equation}
where $\phi(r)$ and $\psi(r)$ are any two functions which are bounded and con\-tinuous in the interval $(a,b)$. By an integration by parts, we have
\begin{equation}\label{eq:6-63}
\int_a^b \phi\,\mathscr{D}_l^2\psi\,r^2\,dr = \bigl[\phi r^2\psi' - \phi'r^2\psi + 2r\phi\psi\bigr]_a^b - \int_a^b \psi\,\mathscr{D}_l^2\phi\,r^2\,dr.
\end{equation}
And by a further integration by parts, we obtain
\begin{equation}\label{eq:6-64}
\int_a^b \phi\,\mathscr{D}_l^4\psi\,r^2\,dr = \bigl[r^2(\phi\mathscr{D}_l^2\psi' - \phi'\mathscr{D}_l^2\psi)\bigr]_a^b + \int_a^b (\mathscr{D}_l^2\phi)(\mathscr{D}_l^2\psi)\,r^2\,dr.
\end{equation}
In particular, if $\phi$ and $\psi$ vanish at both limits,
\begin{equation}\label{eq:6-65}
\int_a^b \phi\,\mathscr{D}_l^2\psi\,r^2\,dr = -\int_a^b \Bigl[\phi'\psi' + \frac{l(l+1)}{r^2}\phi\psi\Bigr]r^2\,dr,
\end{equation}
and
\begin{equation}\label{eq:6-66}
\int_a^b \phi\,\mathscr{D}_l^4\psi\,r^2\,dr = +\int_a^b (\mathscr{D}_l^2\phi)(\mathscr{D}_l^2\psi)\,r^2\,dr.
\end{equation}
This last equation means that \emph{with respect to functions which vanish at the
ends of an interval, the operator $\mathscr{D}_l^4$ is Hermitian}.

Returning to equations~\eqref{eq:6-35} and \eqref{eq:6-36}, let
\begin{equation}\label{eq:6-67}
F = l(l+1)\mathcal{C}\,G,
\end{equation}
the equations then become
\begin{equation}\label{eq:6-68}
\mathscr{D}_l^2(\mathscr{D}_l^2 - \sigma)W = F
\end{equation}
and
\begin{equation}\label{eq:6-69}
(\mathscr{D}_l^2 - p\sigma)F = -l(l+1)\mathcal{C}\,W.
\end{equation}

The boundary conditions are
\begin{equation}\label{eq:6-70}
W = F = 0 \quad\text{and either}\quad\frac{dW}{dr} \quad\text{or}\quad \frac{d^2W}{dr^2} = 0 \quad\text{for}\quad r = 1\quad\text{and}\quad\eta.
\end{equation}
Now multiply equation~\eqref{eq:6-69} by $r^2 F^*$ and integrate over the range
of $r$. Since $F$ vanishes at both limits, we can use equation~\eqref{eq:6-65} and
obtain
\begin{equation}\label{eq:6-71}
\int\Bigl[|F'|^2 + \frac{l(l+1)}{r^2}|F|^2\Bigr]r^2\,dr + p\sigma\int|F|^2 r^2\,dr = l(l+1)\mathcal{C}\int WF^*r^2\,dr.
\end{equation}
By~\eqref{eq:6-68}, the integral on the right-hand side of this equation is
\begin{equation}\label{eq:6-72}
l(l+1)\mathcal{C}\int WF^*r^2\,dr = l(l+1)\mathcal{C}\int W\,\mathscr{D}_l^2(\mathscr{D}_l^2 - \sigma)W^*\,r^2\,dr.
\end{equation}
Transforming the integrals on the right-hand side of equation~\eqref{eq:6-72} with
the aid of equations~\eqref{eq:6-64} and \eqref{eq:6-65}, we obtain
\begin{equation}\label{eq:6-73}
l(l+1)\mathcal{C}\int W F^*\,r^2\,dr = l(l+1)\mathcal{C}\Bigl\{-\Bigl[r^2\Bigl(\frac{dW}{dr}\mathscr{D}_l^2 W^*\Bigr)\Bigr]_a^b + \int|\mathscr{D}_l^2 W|^2 r^2\,dr - \sigma^*\int\Bigl[|W'|^2+\frac{l(l+1)}{r^2}|W|^2\Bigr]r^2\,dr\Bigr\}.
\end{equation}
On a rigid boundary $dW/dr = 0$ and the integrated part in~\eqref{eq:6-73} vanishes;
and if the boundary is free, then $d^2 W/dr^2 = 0$ and
\begin{equation}\label{eq:6-74}
(\mathscr{D}_l^2 W)_{r=\text{const}} = -\frac{l(l+1)}{r^2}W = 0.
\end{equation}
In either case, the integrated part can be written as
\begin{equation}\label{eq:6-75}
\Bigl[r^2\frac{dW}{dr}\frac{d^2W^*}{dr^2}\Bigr].
\end{equation}
Thus equation~\eqref{eq:6-71} reduces to
\begin{equation}\label{eq:6-76}
\int\Bigl[|F'|^2 + \frac{l(l+1)}{r^2}|F|^2\Bigr]r^2\,dr + p\sigma\int|F|^2 r^2\,dr + l(l+1)\mathcal{C}\Bigl\{\int|\mathscr{D}_l^2 W|^2 r^2\,dr - \sigma^*\int\Bigl[|W'|^2+\frac{l(l+1)}{r^2}|W|^2\Bigr]r^2\,dr\Bigr\} = 0.
\end{equation}

The real and the imaginary parts of equation~\eqref{eq:6-76} must vanish separately;
and the vanishing of the imaginary part gives
\begin{equation}\label{eq:6-77}
\sigma_i\Bigl\{p\int|F|^2 r^2\,dr - l(l+1)\mathcal{C}\int\Bigl[|W'|^2+\frac{l(l+1)}{r^2}|W|^2\Bigr]r^2\,dr\Bigr\} = 0.
\end{equation}

The factor of $\sigma_i$ in this equation is positive definite. Accordingly,
\begin{equation}\label{eq:6-78}
\operatorname{Im}(\sigma) = 0.
\end{equation}
Hence $\sigma$ is real and the onset of instability must be via a marginal state
which is stationary. The principle of the exchange of stabilities is
therefore valid.

\section{A variational principle for the case when $\beta$ and $\gamma$ are constants}\label{sec:58}
Since the principle of the exchange of stabilities is valid when $\beta$ and $\gamma$ are
constants, the equations governing the marginal state in this
instance are
\begin{equation}\label{eq:6-79}
\mathscr{D}_l^2 Z = 0,
\end{equation}
and
\begin{equation}\label{eq:6-80}
\mathscr{D}_l^2 F = -l(l+1)\mathcal{C}\,W.
\end{equation}
And the boundary conditions at $r = 1$ and $\eta$ are the same as those given
in equation~\eqref{eq:6-38}.
Equation~\eqref{eq:6-79} clearly requires that
\begin{equation}\label{eq:6-81}
Z = 0.
\end{equation}
The radial component of the vorticity, therefore, vanishes identically.
According to the theorem stated in \S\,\ref{sec:56d}, \emph{the velocity field is purely
poloidal}.

Equations~\eqref{eq:6-60} and \eqref{eq:6-80} together with the boundary conditions~\eqref{eq:6-70}
constitute a characteristic-value problem for $\mathcal{C}$, and instability will set
in with that value of $l$ which leads to the lowest value for $\mathcal{C}$.

The present characteristic-value problem can also be formulated in
terms of a variational principle. The basis for the variational principle
follows from equation~\eqref{eq:6-76} which, when $\sigma = 0$, gives
\begin{equation}\label{eq:6-83}
\mathcal{C} = \frac{\displaystyle\int |\mathscr{D}_l^2 W|^2 r^2\,dr}{\displaystyle l(l+1)\int (W/r)^2\, r^2\,dr}.
\end{equation}
It can be readily proved by the methods described in Chapter~II, \S\,13
that the lowest characteristic value for $\mathcal{C}$ represents the absolute
minimum of the quantity on the right-hand side of equation~\eqref{eq:6-83}. An
equivalent formulation of the principle is that $l(l+1)\mathcal{C}$ is the minimum
value of
\begin{equation}\label{eq:6-84}
\frac{\displaystyle\int |\mathscr{D}_l^2 W|^2 r^2\,dr}{\displaystyle \int |W/r|^2 r^2\,dr},
\end{equation}
for all variations of $F$ which vanish at $r = 1$ and $\eta$ and preserve the
constancy of
\begin{equation}\label{eq:6-85}
\int|W/r|^2 r^2\,dr,
\end{equation}
where $W$ is determined by $F$ and the boundary conditions through the
equation $dW = F$. In this later formulation, $l(l+1)\mathcal{C}$ appears as
the Lagrangian undetermined multiplier.

\subsection{The thermodynamic significance of the variational principle}\label{sec:58a}
We shall now show that the physical content of equation~\eqref{eq:6-83} is the
same as the similar content in earlier chapters: it simply expresses
the equality of the rates at which energy is dissipated by viscosity and is
liberated by the buoyancy force as a sufficient reason for characterizing
the marginal state.

Consider then the rate at which energy is dissipated by viscosity. This
is given by
\begin{equation}\label{eq:6-86}
\epsilon_\nu = \rho\nu\!\int |\operatorname{curl}\boldsymbol{u}|^2\,dV,
\end{equation}
where the integration is extended over the volume of the spherical shell.
We have already seen that the velocity field is entirely poloidal. We
can, therefore, write
\begin{equation}\label{eq:6-87}
\epsilon_\nu = \rho\nu\!\int\!\int |\mathbf{S}|^2 \cdot \cos\theta\,r^2\,dS\,dr.
\end{equation}
By the theorem stated in \S\,\ref{sec:56d}, this is the same as
\begin{equation}\label{eq:6-88}
\epsilon_\nu = \rho\nu\!\int l(l+1)N_l^m\Bigl[\frac{1}{r^2}\Bigl(\frac{d(r\mathscr{S})}{dr}\Bigr)^2 + \frac{l(l+1)}{r^2}\mathscr{S}^2\Bigr] \frac{r^2\,d\theta\,d\varphi}{4\pi}\,dr.
\end{equation}
By an integration by parts, we obtain
\begin{equation}\label{eq:6-89}
\epsilon_\nu = \rho\nu\,l(l+1)N_l^m\!\int \Bigl[\Bigl(\frac{1}{r}\frac{d^2(r\mathscr{S})}{dr^2}\Bigr)^2+\frac{l(l+1)}{r^4}\Bigl(\frac{d(r\mathscr{S})}{dr}\Bigr)^2\Bigr]r^2\,dr.
\end{equation}
Remembering the definition of $\mathscr{S}$ and the requirement that $\mathscr{S}$ vanishes
on the boundaries at $r = R_1$ and $R_0$, we have
\begin{equation}\label{eq:6-90}
\epsilon_\nu = \rho\nu l(l+1)N_l^m\!\int \Bigl[\Bigl(\frac{dW}{dr}\frac{d}{dr}\Bigr)^2 + \frac{l(l+1)}{r^2}W^2\Bigr]\frac{dr}{r^2}.
\end{equation}

From the relation between $W$ and the defining scalar $S$ given in~\eqref{eq:6-58},
it follows that
\begin{equation}\label{eq:6-91}
S = -\frac{l(l+1)}{r}\mathscr{S}\cdot W;
\end{equation}
and we can rewrite equation~\eqref{eq:6-90} in the form
\begin{equation}\label{eq:6-92}
\epsilon_\nu = \frac{\rho\nu\,N_l^m}{l(l+1)}\int |\mathscr{D}_l^2 W|^2\,r^2\,dr.
\end{equation}
Since $W$ and either $dW/dr$ or $d^2W/dr^2$ vanish for $r = 1$ and $r = \eta$, we
can also write
\begin{equation}\label{eq:6-93}
\epsilon_\nu = \frac{\rho\nu N_l^m}{l(l+1)}\int (\mathscr{D}_l^2 W)^2 r^2\,dr.
\end{equation}

Considering next the rate at which internal energy is released by the
buoyancy force acting on the fluid, we must evaluate
\begin{equation}\label{eq:6-94}
\epsilon_T = \alpha\rho g\!\int \Theta u_r\,dV,
\end{equation}
or, according to equation~\eqref{eq:6-11},
\begin{equation}\label{eq:6-95}
\epsilon_T = \alpha\rho\beta\gamma\,N_l^m\!\int \frac{W^2}{r}\,dr,
\end{equation}
where the unit of length is again $R_1$. After performing the angle integra\-tions, we are left with
\begin{equation}\label{eq:6-96}
\epsilon_T = \alpha\rho\beta\gamma\,N_l^m\!\int \frac{W^2}{r^2}\,r^2\,dr.
\end{equation}
Since $G$ vanishes at both limits, we can use equation~\eqref{eq:6-84} to obtain
\begin{equation}\label{eq:6-97}
\epsilon_T = \alpha\rho\beta\gamma N_l^m\!\int |W/r|^2 r^2\,dr.
\end{equation}
Expressing $\mathcal{C}$ in terms of $F$ in accordance with~\eqref{eq:6-61}, we have
\begin{equation}\label{eq:6-98}
\mathcal{C} = \frac{\kappa\nu}{\alpha\beta\gamma R_1^4}.
\end{equation}
Equating $\epsilon_\nu$ and $\epsilon_T$ given by equations~\eqref{eq:6-93} and \eqref{eq:6-97}, we recover equation
\eqref{eq:6-83} which is the basis of the variational principle.

\section{On the onset of thermal instability in a fluid sphere}\label{sec:59}
For a homogeneous fluid sphere, $\beta$ and $\gamma$ are necessarily constants and
the analysis of \S\S\,\ref{sec:57} and \ref{sec:58} applies. We shall now obtain the solution of
the associated characteristic-value problem by the variational method.

Considering the onset of instability as a harmonic of order $l$, we expand
$F$ in a Fourier--Bessel series of the form
\begin{equation}\label{eq:6-99}
F = \sum_{j} A_j J_{l+\frac{1}{2}}(\alpha_j r)/\sqrt{r},
\end{equation}
where $J_{l+\frac{1}{2}}$ denotes the Bessel function of order $l+\frac{1}{2}$ and $\alpha_j$ is its $j$th zero.
The summation over $j$ in equation~\eqref{eq:6-99} may be considered as running from
$1$ to $\infty$; but this is not necessary as we shall confine the $A_j$'s as variational
parameters. With the foregoing choice, the boundary conditions on $F$
at $r = 1$ and $G$ are automatically satisfied.

The functions $J_{l+\frac{1}{2}}(\alpha_j r)$ for $j = 1, 2, \ldots$, and given $l$ form a complete
set of orthonormal functions in the interval $(0, 1)$ and satisfy the ortho\-gonality relation
\begin{equation}\label{eq:6-100}
\int_0^1 J_{l+\frac{1}{2}}(\alpha_j r)J_{l+\frac{1}{2}}(\alpha_k r)\,r\,dr = \delta_{jk}\tfrac{1}{2}[J_{l+\frac{3}{2}}(\alpha_j)]^2,
\end{equation}
where a prime denotes differentiation with respect to the argument.
With the chosen form of $F$, the equation to be solved for $W$ is
\begin{equation}\label{eq:6-101}
\mathscr{D}_l^2 W = -\sum_j A_j J_{l+\frac{1}{2}}(\alpha_j r)/\sqrt{r}.
\end{equation}
We may express the solution of this equation in the form
\begin{equation}\label{eq:6-102}
W = \sum_j A_j W_j,
\end{equation}
where $W_j$ is a solution of the equation
\begin{equation}\label{eq:6-103}
\mathscr{D}_l^2 W_j = -J_{l+\frac{1}{2}}(\alpha_j r)/\sqrt{r},
\end{equation}
which satisfies the necessary boundary conditions. Since
\begin{equation}\label{eq:6-104}
-\frac{J_{l+\frac{1}{2}}(\alpha_j r)}{\sqrt{r}} = -\Bigl(\frac{d^2}{dr^2}+\frac{2}{r}\frac{d}{dr}-\frac{l(l+1)}{r^2}+\alpha_j^2\Bigr)\frac{J_{l+\frac{1}{2}}(\alpha_j r)}{\sqrt{r}},
\end{equation}
a particular integral of equation~\eqref{eq:6-103} which is free of singularity at the
origin is
\begin{equation}\label{eq:6-105}
-\frac{1}{\alpha_j^2}\frac{J_{l+\frac{1}{2}}(\alpha_j r)}{\sqrt{r}}.
\end{equation}
Adding to~\eqref{eq:6-105} the complementary function $B^{(j)} r^l + C^{(j)} r^{-l-1}$ (where
$B^{(j)}$ and $C^{(j)}$ are constants), we have the general solution
\begin{equation}\label{eq:6-106}
W_j = -\frac{1}{\alpha_j^2}\frac{J_{l+\frac{1}{2}}(\alpha_j r)}{\sqrt{r}} + B^{(j)}r^l + C^{(j)}r^{-l-1}.
\end{equation}
The condition $W_j = 0$ at $r = 1$ requires $B^{(j)} = C^{(j)}$ and
\begin{equation}\label{eq:6-107}
\gamma = 2 \quad\text{for a rigid boundary at}\quad r = 1,
\end{equation}
\begin{equation}\label{eq:6-108}
\gamma = 0 \quad\text{for a free boundary at}\quad r = 1.
\end{equation}

We now substitute for $F$ and $W$ according to equations~\eqref{eq:6-99}, \eqref{eq:6-100}, and
\eqref{eq:6-102} in equation~\eqref{eq:6-83}, and obtain
\begin{equation}\label{eq:6-109}
\mathcal{C} = \frac{\displaystyle\sum_j \alpha_j^2 A_j^2 \tfrac{1}{2}[J'_{l+\frac{3}{2}}(\alpha_j)]^2}{\displaystyle l(l+1)\sum_{j,k}A_j A_k\int_0^1 W_j W_k \frac{r^2\,dr}{r^2}}.
\end{equation}
Multiply this equation by $r^2 J_{l+\frac{1}{2}}(\alpha_k r)/\sqrt{r}$ and integrate over the range of $r$.
Making use of the orthogonality relation~\eqref{eq:6-100}, we obtain from equations
\eqref{eq:6-109}, \eqref{eq:6-103}, and \eqref{eq:6-83}, and obtain
\begin{equation}\label{eq:6-111}
\alpha_j^2 \tfrac{1}{2}[J'_{l+\frac{3}{2}}(\alpha_j)]^2 A_j = l(l+1)\mathcal{C}\sum_k Q_{jk} A_k, \quad(j=1,2,\ldots),
\end{equation}
where
\begin{equation}\label{eq:6-112}
Q_{jk} = \int_0^1 W_k \frac{J_{l+\frac{1}{2}}(\alpha_j r)}{\sqrt{r}}\,dr.
\end{equation}

Equation~\eqref{eq:6-111} represents a set of linear homogeneous equations for
the constants $A_j$. This same set of equations ensures that the expression
\eqref{eq:6-83} for $\mathcal{C}$ is a minimum for all variations of the $A_j$'s which preserve the
constancy of~\eqref{eq:6-85}.

Before writing down the secular determinant which follows from equa\-tion~\eqref{eq:6-111}, we shall first evaluate the matrix elements $Q_{jk}$. We shall find
that the integral defining $Q_{jk}$ can be evaluated explicitly if proper use is
made of the various recurrence relations satisfied by the Bessel functions.
We find
\begin{equation}\label{eq:6-113}
B(z) = z\,\frac{J'_{l+\frac{1}{2}}(z)}{J_{l+\frac{1}{2}}(z)}.
\end{equation}

Making further use of the recurrence relations satisfied by the Bessel
functions and remembering that $\alpha_j$ is a zero of $J_{l+\frac{1}{2}}(r)$, we find that
\begin{equation}\label{eq:6-114}
k(z) = -\frac{1}{2l+3}\Bigl(\frac{l+\frac{1}{2}}{z}\Bigr)^2,
\end{equation}
or, substituting for $B^{(j)}$ its value~\eqref{eq:6-106}, we have
\begin{equation}\label{eq:6-115}
k(z) = -\frac{1}{\alpha_j^4}\Bigl[\frac{1}{2}J'^2_{l+\frac{3}{2}}(\alpha_j) + \frac{l+\frac{3}{2}}{\alpha_j^2}J^2_{l+\frac{3}{2}}(\alpha_j)\Bigr],
\end{equation}
which is manifestly symmetric in $k$ and $j$. This symmetry of $Q_{jk}$ reflects
the self-adjoint character of the underlying characteristic-value problem.

With $k(z)$ given by equation~\eqref{eq:6-115}, equation~\eqref{eq:6-111} becomes
\begin{equation}\label{eq:6-117}
\alpha_j^2\tfrac{1}{2}J'^2_{l+\frac{3}{2}}(\alpha_j) = l(l+1)\mathcal{C}\sum_k Q_{jk}\frac{A_k}{A_j}.
\end{equation}

Letting
\begin{equation}\label{eq:6-118}
\epsilon_{jk} = \frac{Q_{jk}}{\tfrac{1}{2}J'^2_{l+\frac{3}{2}}(\alpha_j)\cdot\tfrac{1}{2}J'^2_{l+\frac{3}{2}}(\alpha_k)},
\end{equation}
we can rewrite equation~\eqref{eq:6-117} in the form
\begin{equation}\label{eq:6-119}
\frac{\alpha_j^2}{\tfrac{1}{2}J'^2_{l+\frac{3}{2}}(\alpha_j)} = l(l+1)\mathcal{C}\sum_k \epsilon_{jk}\frac{A_k}{A_j}\cdot\tfrac{1}{2}J'^2_{l+\frac{3}{2}}(\alpha_k).
\end{equation}

The required secular determinant is, therefore,
\begin{equation}\label{eq:6-120}
\bigl|\alpha_j^2\delta_{jk}/[\tfrac{1}{2}J'^2_{l+\frac{3}{2}}(\alpha_j)] - l(l+1)\mathcal{C}\,\epsilon_{jk}\bigr| = 0.
\end{equation}
A first approximation to the value of $\mathcal{C}$ is obtained by setting the $(1,1)$
element of the secular matrix equal to zero. Thus, we obtain
\begin{equation}\label{eq:6-121}
l(l+1)\mathcal{C} = \frac{\alpha_1^2}{\epsilon_{11}}\cdot\frac{1}{[\tfrac{1}{2}J'^2_{l+\frac{3}{2}}(\alpha_1)]}.
\end{equation}
Values in higher approximations can be obtained by retaining more
rows and columns in the secular determinant. The values given in
Tables~XX and XXI were obtained in this manner; those listed under
`exact' in these tables are due to Backus.

It will be observed that in both the cases considered, the easiest modes
to excite are those belonging to $l = 1$.

\subsection{The cell patterns}\label{sec:59a}
The streamlines of the flow are governed by the equations
\begin{equation}\label{eq:6-122}
\frac{dr}{u_r} = \frac{r\,d\theta}{u_\theta} = \frac{r\sin\theta\,d\varphi}{u_\varphi}.
\end{equation}

\begin{table}[ht]
\centering
\caption{The characteristic numbers $\mathcal{C}_l$ and related constants in case the bounding surface is free}
\label{tab:XX}
\begin{tabular}{c c c c c c}
\hline
$l$ & \multicolumn{2}{c}{Natural modes} & Forced & \multicolumn{2}{c}{First approximation} \\
 & First & Second & Third & & \\
\hline
1 & $1.223\times 10^3$ & $1.899\times 10^4$ & $8.162\times 10^4$ & $1.231\times 10^3$ & $1.919\times 10^4$ \\
2 & $1.704\times 10^3$ & $1.988\times 10^4$ & $8.254\times 10^4$ & $1.725\times 10^3$ & $2.010\times 10^4$ \\
3 & $2.432\times 10^3$ & $2.116\times 10^4$ & $8.388\times 10^4$ & $2.470\times 10^3$ & $2.142\times 10^4$ \\
4 & $3.412\times 10^3$ & $2.283\times 10^4$ & $8.565\times 10^4$ & $3.472\times 10^3$ & $2.313\times 10^4$ \\
5 & $4.647\times 10^3$ & $2.487\times 10^4$ & $8.784\times 10^4$ & $4.735\times 10^3$ & $2.523\times 10^4$ \\
6 & $6.139\times 10^3$ & $2.729\times 10^4$ & $9.047\times 10^4$ & $6.261\times 10^3$ & $2.771\times 10^4$ \\
8 & $9.882\times 10^3$ & $3.316\times 10^4$ & $9.702\times 10^4$ & $1.008\times 10^4$ & $3.370\times 10^4$ \\
10 & $1.484\times 10^4$ & $4.042\times 10^4$ & $1.054\times 10^5$ & $1.515\times 10^4$ & $4.111\times 10^4$ \\
12 & $2.102\times 10^4$ & $4.903\times 10^4$ & $1.154\times 10^5$ & $2.147\times 10^4$ & $4.990\times 10^4$ \\
14 & $2.844\times 10^4$ & $5.897\times 10^4$ & $1.270\times 10^5$ & $2.907\times 10^4$ & $6.004\times 10^4$ \\
\hline
\end{tabular}
\end{table}

\begin{table}[ht]
\centering
\caption{The characteristic numbers $\mathcal{C}_l$ and related constants in case the bounding surface is rigid}
\label{tab:XXI}
\begin{tabular}{c c c c c c}
\hline
$l$ & \multicolumn{2}{c}{Natural modes} & Forced & \multicolumn{2}{c}{First approximation} \\
 & First & Second & Third & & \\
\hline
1 & $2.772\times 10^3$ & $3.072\times 10^4$ & $1.172\times 10^5$ & $2.834\times 10^3$ & $3.141\times 10^4$ \\
2 & $3.627\times 10^3$ & $3.180\times 10^4$ & $1.181\times 10^5$ & $3.719\times 10^3$ & $3.255\times 10^4$ \\
3 & $4.855\times 10^3$ & $3.338\times 10^4$ & $1.197\times 10^5$ & $4.989\times 10^3$ & $3.420\times 10^4$ \\
4 & $6.459\times 10^3$ & $3.544\times 10^4$ & $1.217\times 10^5$ & $6.647\times 10^3$ & $3.637\times 10^4$ \\
5 & $8.443\times 10^3$ & $3.799\times 10^4$ & $1.243\times 10^5$ & $8.698\times 10^3$ & $3.904\times 10^4$ \\
6 & $1.081\times 10^4$ & $4.101\times 10^4$ & $1.274\times 10^5$ & $1.115\times 10^4$ & $4.221\times 10^4$ \\
8 & $1.673\times 10^4$ & $4.835\times 10^4$ & $1.352\times 10^5$ & $1.730\times 10^4$ & $4.985\times 10^4$ \\
10 & $2.428\times 10^4$ & $5.736\times 10^4$ & $1.450\times 10^5$ & $2.515\times 10^4$ & $5.923\times 10^4$ \\
12 & $3.347\times 10^4$ & $6.803\times 10^4$ & $1.568\times 10^5$ & $3.471\times 10^4$ & $7.032\times 10^4$ \\
14 & $4.431\times 10^4$ & $8.033\times 10^4$ & $1.705\times 10^5$ & $4.602\times 10^4$ & $8.311\times 10^4$ \\
\hline
\end{tabular}
\end{table}

\figurePlaceholder{56}{The Rayleigh number $\mathcal{C}_l$ for the onset of thermal convection in a fluid sphere as a spherical harmonic disturbance of order $l$ for the two cases: when the bounding surface is rigid (curve labelled $a$) and when the bounding surface is free (curve labelled $b$).}

When the velocity field is purely poloidal (as in the present instance),
the equations take the explicit forms
\begin{equation}\label{eq:6-123}
u_r = l(l+1)\frac{\mathscr{S}}{r}Y_l^m(\theta,\varphi),\quad u_\theta = \frac{1}{r}\frac{d(r\mathscr{S})}{dr}\frac{\partial Y_l^m}{\partial\theta},\quad u_\varphi = \frac{im}{r\sin\theta}\frac{d(r\mathscr{S})}{dr}Y_l^m.
\end{equation}
These equations can be explicitly integrated for the lowest axisymmetric

\figurePlaceholder{57}{The pattern of convection in a fluid sphere for the lowest symmetric mode of disturbance, $l=1$, $m=0$: (a) for a sphere with a rigid boundary, and (b) for a sphere with a free boundary.}

mode $l = 1$ and $m = 0$: the streamlines are then confined to the meri\-dional planes; and in these planes the equation governing them is
\begin{equation}\label{eq:6-124}
\frac{1}{r^2}\frac{d\mathscr{S}}{dr}\sin^2\theta = -\text{const}\cdot P_1\sin\theta\,d\theta.
\end{equation}
The integral of this equation is
\begin{equation}\label{eq:6-125}
r\mathscr{S}\sin\theta = \text{constant.}
\end{equation}
In terms of $W$ the equation is
\begin{equation}\label{eq:6-126}
W\sin^2\theta/r = \text{constant.}
\end{equation}
The streamlines derived from this equation are illustrated in Fig.~57\,$a$, $b$.

\section{On the onset of thermal instability in spherical shells}\label{sec:60}
On the assumption that the principle of the exchange of stabilities is
valid, the equations governing the marginal states are (cf.\ equations~\eqref{eq:6-35}
and \eqref{eq:6-36}):
\begin{equation}\label{eq:6-127}
\mathscr{D}_l^2\,G = -\frac{\beta}{r}\,W
\end{equation}
and
\begin{equation}\label{eq:6-128}
\mathscr{D}_l^2 W = l(l+1)\frac{\psi}{\nu}\,G,
\end{equation}
where $R_0$ is the radius of the outer spherical boundary. Letting
\begin{equation}\label{eq:6-129}
\beta(r) = \bar{\beta}_0\, b(r),\quad \gamma(r) = \gamma_0\, c(r),
\end{equation}
and
\begin{equation}\label{eq:6-130}
F = b(r)\,l(l+1)\mathcal{C}\,G,
\end{equation}
where $\bar{\beta}_0$ and $\gamma_0$ are the values of $\beta(r)$ and $\gamma(r)$ at $r = 1$, we can rewrite
equations~\eqref{eq:6-127} and \eqref{eq:6-128} in the forms
\begin{equation}\label{eq:6-131}
\mathscr{D}_l^2 W = c(r)\,F,
\end{equation}
and
\begin{equation}\label{eq:6-132}
\mathscr{D}_l^2 F = -b(r)\,l(l+1)\mathcal{C}\,W,
\end{equation}
where
\begin{equation}\label{eq:6-133}
\mathcal{C} = \frac{\alpha\bar{\beta}_0\gamma_0 R_0^4}{\kappa\nu}.
\end{equation}

Solutions of equations~\eqref{eq:6-131} and \eqref{eq:6-132} must be sought which satisfy the
boundary conditions~\eqref{eq:6-70}.

The form of equation~\eqref{eq:6-131} which must be solved for $W$ suggests that
we expand $F$ in a series of cylinder functions of order $l+\frac{1}{2}$ which vanish
at $r = 1$ and $\eta$. Such functions can be constructed as follows. Let
\begin{equation}\label{eq:6-134}
\mathscr{W}_{l+\frac{1}{2}}(r,\alpha) = J_{l+\frac{1}{2}}(\alpha r)Y_{l+\frac{1}{2}}(\alpha) - J_{l+\frac{1}{2}}(\alpha)Y_{l+\frac{1}{2}}(\alpha r),
\end{equation}
where $\alpha$ is a constant unspecified for the present. Then,
\begin{equation}\label{eq:6-135}
\mathscr{W}_{l+\frac{1}{2}}(1,\alpha) = J_{l+\frac{1}{2}}(\alpha)Y_{l+\frac{1}{2}}(\alpha) - J_{l+\frac{1}{2}}(\alpha)Y_{l+\frac{1}{2}}(\alpha) = 0,
\end{equation}
clearly vanishes at $r = 1$; and it will also vanish at $r = \eta > 0$
\begin{equation}\label{eq:6-136}
\mathscr{W}_{l+\frac{1}{2}}(\eta,\alpha) = 0.
\end{equation}

It is known that equation~\eqref{eq:6-136} admits an infinity of roots, all of which
are real and simple; and further, that if $\alpha_j$ ($j = 1, 2, \ldots$) are the distinct
roots of the equation, the functions $\mathscr{W}_{l+\frac{1}{2}}(\alpha_j r)/\sqrt{r}$ form an orthogonal set
of functions with the integral property
\begin{equation}\label{eq:6-137}
\int_\eta^1 \frac{\mathscr{W}_{l+\frac{1}{2}}(\alpha_j r)}{\sqrt{r}}\cdot\frac{\mathscr{W}_{l+\frac{1}{2}}(\alpha_k r)}{\sqrt{r}}\,r^2\,dr = N_{jk},
\end{equation}
where
\begin{equation}\label{eq:6-138}
N_{jk} = \delta_{jk}\cdot\frac{2}{\pi\alpha_j^2}\Bigl[\frac{J^2_{l+\frac{1}{2}}(\alpha_j\eta)}{J^2_{l+\frac{1}{2}}(\alpha_j)} - 1\Bigr].
\end{equation}
For later use, we may note here that the derivatives of $\mathscr{W}_{l+\frac{1}{2}}(\alpha_j r)$
at $r = 1$ and $r = \eta$ (which we shall denote by $\mathscr{W}'_{l+\frac{1}{2}}(1)$ and $\mathscr{W}'_{l+\frac{1}{2}}(\eta)$,
respectively) are given by
\begin{equation}\label{eq:6-139}
\mathscr{W}'_{l+\frac{1}{2}}(1) = -\frac{2}{\pi},
\end{equation}
and
\begin{equation}\label{eq:6-140}
\mathscr{W}'_{l+\frac{1}{2}}(\eta) = \frac{2}{\pi\eta}\frac{J_{l+\frac{1}{2}}(\alpha_j)}{J_{l+\frac{1}{2}}(\alpha_j\eta)}.
\end{equation}
Also, since $\mathscr{W}_{l+\frac{1}{2}}$ represents a solution of Bessel's equation of order
$\nu$, it must clearly satisfy the usual recurrence relations implied by the
equation
\begin{equation}\label{eq:6-141}
\frac{d}{dr}\Bigl[\frac{1}{r}\mathscr{W}_{l+\frac{1}{2}}\Bigr] = -\frac{\alpha}{r}\mathscr{W}_{l+\frac{3}{2}},
\end{equation}
and
\begin{equation}\label{eq:6-142}
\frac{d}{dr}\Bigl[r\,\mathscr{W}_{l+\frac{1}{2}}\Bigr] = \alpha r\,\mathscr{W}_{l-\frac{1}{2}}.
\end{equation}

Returning to the solution of equations~\eqref{eq:6-131} and \eqref{eq:6-132}, we express $F$
as a series in the form
\begin{equation}\label{eq:6-143}
F = \sum_j A_j\frac{\mathscr{W}_{l+\frac{1}{2}}(\alpha_j r)}{\sqrt{r}},
\end{equation}
where the $A_j$'s are constants. With this choice for $F$, the solution for $G$
can be written in the form
\begin{equation}\label{eq:6-144}
W = \sum_j A_j W_j,
\end{equation}
where $W_j$ is a solution of the equation
\begin{equation}\label{eq:6-145}
\mathscr{D}_l^2 W_j = -\frac{1}{\alpha_j^2}\frac{\mathscr{W}_{l+\frac{1}{2}}(\alpha_j r)}{\sqrt{r}};
\end{equation}
the general solution of equation~\eqref{eq:6-144} is given by
\begin{equation}\label{eq:6-146}
W_j = -\frac{1}{\alpha_j^2}\frac{\mathscr{W}_{l+\frac{1}{2}}(\alpha_j r)}{\sqrt{r}} + B_j^{(1)} r^l + B_j^{(2)} r^{-l-1},
\end{equation}
where $B_j^{(1)}$, $B_j^{(2)}$ are constants of integration to be determined by the
boundary conditions. The conditions $W = 0$ at $r = 1$ and $r = \eta$
(which do not depend on the nature of the bounding surfaces) require
\begin{equation}\label{eq:6-147}
B_j^{(1)} + B_j^{(2)} = 0
\end{equation}
and
\begin{equation}\label{eq:6-148}
B_j^{(1)}\eta^l + B_j^{(2)}\eta^{-l-1} = 0.
\end{equation}
The remaining boundary conditions depend on the nature of the surfaces
at $r = 1$ and $r = \eta$; and there are four cases to consider. We shall return
to them presently.

We now substitute for $F$ and $W$ in accordance with
equations~\eqref{eq:6-143}, \eqref{eq:6-144}, and \eqref{eq:6-146} into equation~\eqref{eq:6-132} and obtain

\begin{equation}\label{eq:6-149}
\alpha_j^2 N_{jj} A_j = l(l+1)\mathcal{C}\sum_k Q_{jk} A_k,
\end{equation}

Multiply this equation by $r^2\,\mathscr{W}_{l+\frac{1}{2}}(\alpha_j r)/\sqrt{r}$ and integrate over the range of $r$.
We obtain
\begin{equation}\label{eq:6-150}
\alpha_j^2 N_{jj} A_j = l(l+1)\mathcal{C}\sum_k Q_{jk} A_k,
\end{equation}
where
\begin{equation}\label{eq:6-151}
Q_{jk} = \int_\eta^1 W_k\,\frac{\mathscr{W}_{l+\frac{1}{2}}(\alpha_j r)}{\sqrt{r}}\,dr.
\end{equation}
Letting
\begin{equation}\label{eq:6-152}
P_{jk} = -\int_\eta^1\frac{\mathscr{W}_{l+\frac{1}{2}}(\alpha_j r)}{\sqrt{r}}\frac{\mathscr{W}_{l+\frac{1}{2}}(\alpha_k r)}{\sqrt{r}}\,\frac{dr}{r},
\end{equation}
we can rewrite equation~\eqref{eq:6-150} in the form
\begin{equation}\label{eq:6-153}
\sum_k\Bigl[P_{jk}-\alpha_k^2 N_{jj}\delta_{jk}+l(l+1)\mathcal{C}\,Q_{jk}\Bigr]A_k = 0.
\end{equation}

This leads to the secular equation
\begin{equation}\label{eq:6-154}
\bigl|P_{jk}-\alpha_k^2 N_{jj}\delta_{jk}+l(l+1)\mathcal{C}\,Q_{jk}\bigr| = 0.
\end{equation}

From the Hermitian character of the operator $\mathscr{D}_l^2$ with respect to
functions which vanish at the ends of an interval, it follows that
\begin{equation}\label{eq:6-155}
P_{jk} = -S_{jk}\frac{1}{\alpha_j^2\alpha_k^2}\bigl[J'_{l+\frac{1}{2}}(\alpha_j)\bigr]^2N_{jk}^2.
\end{equation}
Hence,
\begin{equation}\label{eq:6-156}
P_{jk} = P_{kj}.
\end{equation}

The matrix element $Q_{jk}$ can be explicitly evaluated if proper use is
made of the recurrence relations satisfied by the cylinder functions. We
find after some lengthy but straightforward reductions that
\begin{equation}\label{eq:6-157}
Q_{jk} = \Bigl[\frac{\alpha_k}{\alpha_j^2-\alpha_k^2}\bigl(r^{-\frac{1}{2}}\mathscr{W}_{l+\frac{1}{2}}(\alpha_j r)\mathscr{W}_{l+\frac{3}{2}}(\alpha_k r) + \cdots\bigr)\Bigr]_\eta^1\frac{1}{\alpha_k^2} N_{jk}^2,
\end{equation}
where the boundary conditions leading to equations~\eqref{eq:6-147} and \eqref{eq:6-148} have
contributed to the simplification.

From the recurrence relations satisfied by the functions $\mathscr{W}_{l+\frac{1}{2}}(r)$ and
the fact that $\mathscr{W}_{l+\frac{1}{2}}(1)$ and $\mathscr{W}_{l+\frac{1}{2}}(\eta)$ are zero, it follows that
\begin{equation}\label{eq:6-158}
\mathscr{W}'_{l+\frac{3}{2}}(1) = -(2l+3)\mathscr{W}_{l+\frac{3}{2}}(1),
\end{equation}
and
\begin{equation}\label{eq:6-159}
\mathscr{W}'_{l+\frac{3}{2}}(\eta) = \frac{2l+3}{\eta}\mathscr{W}_{l+\frac{3}{2}}(\eta).
\end{equation}
Further, let
\begin{equation}\label{eq:6-160}
\mathscr{K}_{jk} = \mathscr{W}_{l+\frac{3}{2}}(1,\alpha_j)\mathscr{W}_{l+\frac{3}{2}}(1,\alpha_k) - (2l+3)\mathscr{W}_{l+\frac{3}{2}}(\eta,\alpha_j)\mathscr{W}_{l+\frac{3}{2}}(\eta,\alpha_k),
\end{equation}
and
\begin{equation}\label{eq:6-161}
Q_{jk} = \frac{1}{2}(2l+3)\bigl[\mathscr{K}_{jk}\alpha_j\alpha_k R_0^{-2} - (2l-1)\mathscr{W}_{l+\frac{1}{2}}\cdot\mathscr{W}_{l+\frac{1}{2}} N_{jk}\bigr]\frac{1}{\alpha_k^2}.
\end{equation}

The further reduction of $Q_{jk}$ requires an explicit consideration of the
nature of the surfaces at $r = 1$ and $\eta$; and we consider separately the
four possible cases.

\medskip
\noindent (i) \emph{Free surfaces at $r = 1$ and $r = \eta$}

In this case, $d^2 W/dr^2 = 0$ at $r = 1$ and $r = \eta$; and the conditions
applied to the solution~\eqref{eq:6-146} give
\begin{equation}\label{eq:6-162}
B_j^{(1)}l(l-1) + B_j^{(2)}(l+1)(l+2) = \frac{1}{\alpha_j^2}\frac{\mathscr{W}''_{l+\frac{1}{2}}(1)}{\sqrt{1}}
\end{equation}
and
\begin{equation}\label{eq:6-163}
B_j^{(1)}l(l-1)\eta^{l-2} + B_j^{(2)}(l+1)(l+2)\eta^{-l-3} = 2\gamma_0 + \frac{\mathscr{W}''_{l+\frac{1}{2}}(\alpha_j\eta)}{\sqrt{\eta}\,\alpha_j^2}.
\end{equation}
Solving these equations together with equations~\eqref{eq:6-147} and \eqref{eq:6-148}, we find
\begin{equation}\label{eq:6-164}
B_j^{(1)} = \frac{\mathscr{W}''_{l+\frac{1}{2}}(1)}{(2l+1)\alpha_j^2},
\end{equation}
\begin{equation}\label{eq:6-165}
B_j^{(2)} = \frac{-\mathscr{W}''_{l+\frac{1}{2}}(\eta)\eta^{l+1}}{(2l+1)\alpha_j^2}.
\end{equation}
Inserting the foregoing expressions for $B_j^{(1)}$ and $B_j^{(2)}$ in equation~\eqref{eq:6-161} for
$Q_{jk}$, we obtain
\begin{equation}\label{eq:6-166}
Q_{jk} = \text{(explicit expression symmetric in } j \text{ and } k\text{)}.
\end{equation}
We observe that $Q_{jk}$ is symmetric in $k$ and $j$.

\medskip
\noindent (ii) \emph{A rigid surface at $r = \eta$ and a free surface at $r = 1$}

In this case $dW/dr = 0$ at $r = \eta$ and $d^2W/dr^2 = 0$ at $r = 1$; and these
conditions applied to the solution~\eqref{eq:6-146} give
\begin{equation}\label{eq:6-167}
B_j^{(1)} l\eta^{l-1} - B_j^{(2)}(l+1)\eta^{-l-2} = \frac{1}{\alpha_j^2}\frac{d}{dr}\Bigl[\frac{\mathscr{W}_{l+\frac{1}{2}}(\alpha_j r)}{\sqrt{r}}\Bigr]_{r=\eta}
\end{equation}
and
\begin{equation}\label{eq:6-168}
l(l-1)B_j^{(1)} + (l+1)(l+2)B_j^{(2)} = -(l+1)\frac{\mathscr{W}''_{l+\frac{1}{2}}(1)}{\alpha_j^2}.
\end{equation}
Solving these equations together with equations~\eqref{eq:6-147} and \eqref{eq:6-148}, we find
\begin{equation}\label{eq:6-169}
B_j^{(1)} = \text{(expression)},
\end{equation}
\begin{equation}\label{eq:6-170}
B_j^{(2)} = \text{(expression)},
\end{equation}
where
\begin{equation}\label{eq:6-171}
\Delta_l(r,\eta) = l(l+1)r^{l-1}\eta^{-l-2} + (l+1)\eta^{l-1}r^{-l-2}.
\end{equation}
Inserting the foregoing expressions for $B_j^{(1)}$ and $B_j^{(2)}$ in equation~\eqref{eq:6-161} for
$Q_{jk}$, we obtain (cf.\ equation~\eqref{eq:6-175}):
\begin{equation}\label{eq:6-172}
Q_{jk} = -\frac{1}{\alpha_j^2\alpha_k^2}\Bigl[\bigl(l+1\bigr)(3l+1)\bigl(\frac{l+\frac{1}{2}}{\eta}\bigr) \cdots\Bigr],
\end{equation}
which is clearly symmetric in $k$ and $j$.

\medskip
\noindent (iii) \emph{A free surface at $r = \eta$ and a rigid surface at $r = 1$}

In this case $d^2 W/dr^2 = 0$ at $r = \eta$ and $dW/dr = 0$ at $r = 1$; and these
conditions applied to the solution~\eqref{eq:6-146} give
\begin{equation}\label{eq:6-173}
l B_j^{(1)} - (l+1)B_j^{(2)} = \frac{1}{\alpha_j^2}\frac{d}{dr}\Bigl[\frac{\mathscr{W}_{l+\frac{1}{2}}(\alpha_j r)}{\sqrt{r}}\Bigr]_{r=1}
\end{equation}
and
\begin{equation}\label{eq:6-174}
l(l-1)\eta^{l-2}B_j^{(1)} + (l+1)(l+2)\eta^{-l-3}B_j^{(2)} = -(l+1)\frac{\mathscr{W}''_{l+\frac{1}{2}}(\alpha_j\eta)}{\sqrt{\eta}\,\alpha_j^2}.
\end{equation}
Solving these equations together with equations~\eqref{eq:6-147} and \eqref{eq:6-148}, we find
\begin{equation}\label{eq:6-175}
B_j^{(1)} = \text{(expression)},
\end{equation}
\begin{equation}\label{eq:6-176}
B_j^{(2)} = \text{(expression)},
\end{equation}
where $\Delta_l(r,\eta)$ has the same meaning as in equation~\eqref{eq:6-171}. Inserting the
foregoing expressions~\eqref{eq:6-175} for $B_j^{(1)}$ and $B_j^{(2)}$ in equation~\eqref{eq:6-161} for $Q_{jk}$, we obtain
(cf.\ equation~\eqref{eq:6-172}):
\begin{equation}\label{eq:6-177}
Q_{jk} = \text{(expression symmetric in $j$ and $k$)}.
\end{equation}

\medskip
\noindent (iv) \emph{Rigid surfaces at $r = 1$ and $r = \eta$}

In this case $dW/dr = 0$ at $r = 1$ and $r = \eta$; and these conditions
applied to the solution~\eqref{eq:6-146} give
\begin{equation}\label{eq:6-178}
l B_j^{(1)} - (l+1) B_j^{(2)} = \frac{1}{\alpha_j^2}\frac{d}{dr}\Bigl[\frac{\mathscr{W}_{l+\frac{1}{2}}(\alpha_j r)}{\sqrt{r}}\Bigr]_{r=1}
\end{equation}
and
\begin{equation}\label{eq:6-179}
l\eta^{l-1}B_j^{(1)} - (l+1)\eta^{-l-2}B_j^{(2)} = \frac{1}{\alpha_j^2}\frac{d}{dr}\Bigl[\frac{\mathscr{W}_{l+\frac{1}{2}}(\alpha_j r)}{\sqrt{r}}\Bigr]_{r=\eta}.
\end{equation}
Solving these equations together with equations~\eqref{eq:6-147} and \eqref{eq:6-148}, we find
\begin{equation}\label{eq:6-180}
B_j^{(1)} = \text{(expression)},
\end{equation}
and
\begin{equation}\label{eq:6-181}
B_j^{(2)} = \text{(expression)},
\end{equation}
where
\begin{equation}\label{eq:6-182}
\Delta_0(r,\eta) = (4l^2+4l+3)\bigl(r^{l-1}\eta^{-l-2}+r^{-l-2}\eta^{l-1}\bigr),
\end{equation}
Inserting the foregoing expressions for $B_j^{(1)}$ and $B_j^{(2)}$ in equation~\eqref{eq:6-161} for
$Q_{jk}$, we obtain:
\begin{equation}\label{eq:6-183}
Q_{jk} = \text{(expression)},
\end{equation}
which is again symmetric in $k$ and $j$.

\subsection{The case $b = c = 1$}\label{sec:60a}
In this case we know that the underlying characteristic-value problem
is self-adjoint. Consequently, solving the secular equation~\eqref{eq:6-154} is
equivalent to regarding the $A_j$'s as variational parameters and mini\-mizing the expression~\eqref{eq:6-83} for $\mathcal{C}$.

\figurePlaceholder{58}{The Rayleigh number $\mathcal{C}_l$ for the onset of convection in the mantle as a spherical harmonic disturbance of order $l$ for various thicknesses of the mantle. The different curves are distinguished by the values of the fraction ($\eta$) of the radius of the sphere occupied by the core. The following cases are distinguished: (a) free surfaces at $r = 1$ and $r = \eta$; (b) a rigid surface at $r = \eta$ and a free surface at $r = 1$; (c) a free surface at $r = \eta$ and a rigid surface at $r = 1$; (d) rigid surfaces at $r = 1$ and $r = \eta$.}

When $b = c = 1$,
\begin{equation}\label{eq:6-184}
P_{jk} = \int_\eta^1 \frac{\mathscr{W}_{l+\frac{1}{2}}(\alpha_j r)}{\sqrt{r}}\frac{\mathscr{W}_{l+\frac{1}{2}}(\alpha_k r)}{\sqrt{r}}\frac{dr}{r^2} = \delta_{jk} N_{jk}\alpha_j^2.
\end{equation}
The characteristic equation~\eqref{eq:6-154} can in this case be rewritten in the form
\begin{equation}\label{eq:6-185}
\Bigl|\frac{\delta_{jk}}{\alpha_k^2 N_{kk}} - \frac{l(l+1)\mathcal{C}}{P_{jk}}\cdot Q_{jk}\Bigr| = 0.
\end{equation}

The symmetry of this matrix is a reflection of the self-adjoint character
of the problem we are considering.

\begin{table}[ht]
\centering
\caption{The characteristic numbers $\mathcal{C}_l$ for various values of $l$ and $\eta$\newline
(Free surfaces at $r = 1$ and $r = \eta$)}
\label{tab:XXII}
\begin{tabular}{c c c c c c}
\hline
$l$ & $\eta = 0$ & $\eta = 0.1$ & $\eta = 0.3$ & $\eta = 0.5$ & $\eta = 0.8$ \\
\hline
1 & $1.223\times 10^3$ & $1.31\times 10^3$ & $2.06\times 10^3$ & $6.12\times 10^3$ & $1.24\times 10^5$ \\
2 & $1.704\times 10^3$ & $1.70\times 10^3$ & $1.58\times 10^3$ & $2.60\times 10^3$ & $2.24\times 10^4$ \\
3 & $2.432\times 10^3$ & $2.33\times 10^3$ & $1.65\times 10^3$ & $1.73\times 10^3$ & $7.52\times 10^3$ \\
4 & $3.412\times 10^3$ & $3.21\times 10^3$ & $1.91\times 10^3$ & $1.43\times 10^3$ & $3.86\times 10^3$ \\
5 & $4.647\times 10^3$ & $4.33\times 10^3$ & $2.29\times 10^3$ & $1.32\times 10^3$ & $2.49\times 10^3$ \\
6 & $6.139\times 10^3$ & $5.68\times 10^3$ & $2.77\times 10^3$ & $1.31\times 10^3$ & $1.85\times 10^3$ \\
8 & $9.882\times 10^3$ & $9.08\times 10^3$ & $4.01\times 10^3$ & $1.46\times 10^3$ & $1.26\times 10^3$ \\
10 & $1.484\times 10^4$ & $1.36\times 10^4$ & $5.66\times 10^3$ & $1.80\times 10^3$ & $1.09\times 10^3$ \\
12 & $2.102\times 10^4$ & $1.92\times 10^4$ & $7.70\times 10^3$ & $2.27\times 10^3$ & $1.04\times 10^3$ \\
14 & $2.844\times 10^4$ & $2.59\times 10^4$ & $1.01\times 10^4$ & $2.84\times 10^3$ & $1.05\times 10^3$ \\
\hline
\end{tabular}
\end{table}

In the first approximation, equation~\eqref{eq:6-185} gives
\begin{equation}\label{eq:6-186}
l(l+1)\mathcal{C} = \frac{\alpha_1^4 N_{11}}{Q_{11}}.
\end{equation}

This formula, together with the expressions for $Q_{jk}$ given earlier, has
been used to determine the characteristic numbers $\mathcal{C}_l$ for various values
of $l$ and $\eta$ and for the sets of boundary conditions we have considered.
The results are given in Tables~XXII--XXV; they are further illustrated
in Fig.~58\,$a$--$d$. It is apparent that as the thickness of the shell
decreases, the pattern of the convection which manifests itself at
marginal stability shifts progressively to harmonics of the higher orders.
As we have stated, the results given in Tables~XXII--XXV have been
obtained only in the first approximation. Nevertheless, we may expect
from past experience that the values tabulated are not in error by more
than 2 or 3 per cent; this is confirmed by the comparison made in
Table~XXVI of the values obtained in the first and the second approxi\-mations for case~(i) and $\eta = 0.5$.

\begin{table}[ht]
\centering
\caption{The characteristic numbers $\mathcal{C}_l$ for various values of $l$ and $\eta$\newline
(A rigid surface at $r = \eta$ and a free surface at $r = 1$)}
\label{tab:XXIII}
\begin{tabular}{c c c c c c}
\hline
$l$ & $\eta = 0$ & $\eta = 0.1$ & $\eta = 0.3$ & $\eta = 0.5$ & $\eta = 0.8$ \\
\hline
1 & $1.223\times 10^3$ & $1.36\times 10^3$ & $2.88\times 10^3$ & $1.21\times 10^4$ & $3.68\times 10^5$ \\
2 & $1.704\times 10^3$ & $1.81\times 10^3$ & $2.28\times 10^3$ & $5.22\times 10^3$ & $6.64\times 10^4$ \\
3 & $2.432\times 10^3$ & $2.51\times 10^3$ & $2.22\times 10^3$ & $3.25\times 10^3$ & $2.08\times 10^4$ \\
4 & $3.412\times 10^3$ & $3.47\times 10^3$ & $2.44\times 10^3$ & $2.51\times 10^3$ & $9.98\times 10^3$ \\
5 & $4.647\times 10^3$ & $4.69\times 10^3$ & $2.83\times 10^3$ & $2.19\times 10^3$ & $6.05\times 10^3$ \\
6 & $6.139\times 10^3$ & $6.16\times 10^3$ & $3.34\times 10^3$ & $2.06\times 10^3$ & $4.22\times 10^3$ \\
8 & $9.882\times 10^3$ & $9.82\times 10^3$ & $4.68\times 10^3$ & $2.08\times 10^3$ & $2.67\times 10^3$ \\
10 & $1.484\times 10^4$ & $1.47\times 10^4$ & $6.41\times 10^3$ & $2.33\times 10^3$ & $2.07\times 10^3$ \\
12 & $2.102\times 10^4$ & $2.08\times 10^4$ & $8.53\times 10^3$ & $2.73\times 10^3$ & $1.81\times 10^3$ \\
14 & $2.844\times 10^4$ & $2.81\times 10^4$ & $1.10\times 10^4$ & $3.25\times 10^3$ & $1.72\times 10^3$ \\
\hline
\end{tabular}
\end{table}

\begin{table}[ht]
\centering
\caption{The characteristic numbers $\mathcal{C}_l$ for various values of $l$ and $\eta$\newline
(A free surface at $r = \eta$ and a rigid surface at $r = 1$)}
\label{tab:XXIV}
\begin{tabular}{c c c c c c}
\hline
$l$ & $\eta = 0$ & $\eta = 0.1$ & $\eta = 0.3$ & $\eta = 0.5$ & $\eta = 0.8$ \\
\hline
1 & $2.772\times 10^3$ & $2.86\times 10^3$ & $3.56\times 10^3$ & $9.48\times 10^3$ & $1.69\times 10^5$ \\
2 & $3.627\times 10^3$ & $3.55\times 10^3$ & $2.80\times 10^3$ & $4.08\times 10^3$ & $3.08\times 10^4$ \\
3 & $4.855\times 10^3$ & $4.64\times 10^3$ & $2.78\times 10^3$ & $2.72\times 10^3$ & $1.06\times 10^4$ \\
4 & $6.459\times 10^3$ & $6.09\times 10^3$ & $3.06\times 10^3$ & $2.22\times 10^3$ & $5.42\times 10^3$ \\
5 & $8.443\times 10^3$ & $7.90\times 10^3$ & $3.51\times 10^3$ & $2.01\times 10^3$ & $3.49\times 10^3$ \\
6 & $1.081\times 10^4$ & $1.01\times 10^4$ & $4.09\times 10^3$ & $1.94\times 10^3$ & $2.60\times 10^3$ \\
8 & $1.673\times 10^4$ & $1.55\times 10^4$ & $5.59\times 10^3$ & $2.05\times 10^3$ & $1.79\times 10^3$ \\
10 & $2.428\times 10^4$ & $2.24\times 10^4$ & $7.53\times 10^3$ & $2.38\times 10^3$ & $1.51\times 10^3$ \\
12 & $3.347\times 10^4$ & $3.08\times 10^4$ & $9.90\times 10^3$ & $2.85\times 10^3$ & $1.41\times 10^3$ \\
14 & $4.431\times 10^4$ & $4.08\times 10^4$ & $1.27\times 10^4$ & $3.44\times 10^3$ & $1.39\times 10^3$ \\
\hline
\end{tabular}
\end{table}

\begin{table}[ht]
\centering
\caption{The characteristic numbers $\mathcal{C}_l$ for various values of $l$ and $\eta$\newline
(Rigid surfaces at $r = 1$ and $r = \eta$)}
\label{tab:XXV}
\begin{tabular}{c c c c c c}
\hline
$l$ & $\eta = 0$ & $\eta = 0.1$ & $\eta = 0.3$ & $\eta = 0.5$ & $\eta = 0.8$ \\
\hline
1 & $2.772\times 10^3$ & $2.95\times 10^3$ & $4.71\times 10^3$ & $1.81\times 10^4$ & $4.93\times 10^5$ \\
2 & $3.627\times 10^3$ & $3.73\times 10^3$ & $3.72\times 10^3$ & $7.75\times 10^3$ & $8.97\times 10^4$ \\
3 & $4.855\times 10^3$ & $4.89\times 10^3$ & $3.56\times 10^3$ & $4.80\times 10^3$ & $2.85\times 10^4$ \\
4 & $6.459\times 10^3$ & $6.44\times 10^3$ & $3.77\times 10^3$ & $3.63\times 10^3$ & $1.38\times 10^4$ \\
5 & $8.443\times 10^3$ & $8.38\times 10^3$ & $4.18\times 10^3$ & $3.10\times 10^3$ & $8.35\times 10^3$ \\
6 & $1.081\times 10^4$ & $1.07\times 10^4$ & $4.76\times 10^3$ & $2.85\times 10^3$ & $5.85\times 10^3$ \\
8 & $1.673\times 10^4$ & $1.65\times 10^4$ & $6.31\times 10^3$ & $2.74\times 10^3$ & $3.66\times 10^3$ \\
10 & $2.428\times 10^4$ & $2.39\times 10^4$ & $8.32\times 10^3$ & $2.95\times 10^3$ & $2.82\times 10^3$ \\
12 & $3.347\times 10^4$ & $3.30\times 10^4$ & $1.08\times 10^4$ & $3.33\times 10^3$ & $2.43\times 10^3$ \\
14 & $4.431\times 10^4$ & $4.37\times 10^4$ & $1.37\times 10^4$ & $3.85\times 10^3$ & $2.26\times 10^3$ \\
\hline
\end{tabular}
\end{table}

\begin{table}[ht]
\centering
\caption{The characteristic numbers $\mathcal{C}_l$ in the first and the second approximations
for various values of $l$ and $\eta = 0.5$\newline
(Free surfaces at $r = 1$ and $r = \eta$)}
\label{tab:XXVI}
\begin{tabular}{c c c}
\hline
$l$ & First approximation & Second approximation \\
\hline
1 & $6.12\times 10^3$ & $6.05\times 10^3$ \\
2 & $2.60\times 10^3$ & $2.58\times 10^3$ \\
3 & $1.73\times 10^3$ & $1.72\times 10^3$ \\
4 & $1.43\times 10^3$ & $1.42\times 10^3$ \\
5 & $1.32\times 10^3$ & $1.31\times 10^3$ \\
6 & $1.31\times 10^3$ & $1.30\times 10^3$ \\
8 & $1.46\times 10^3$ & $1.45\times 10^3$ \\
10 & $1.80\times 10^3$ & $1.78\times 10^3$ \\
\hline
\end{tabular}
\end{table}

\subsection{The case $b(r) = 1$}\label{sec:60b}
This case corresponds to a uniform distribution of the heat sources
and the allowed non-constancy of $c(r)$ permits a variation of gravity
in the mantle which is different from that in a homogeneous sphere.

When $b = 1$, the matrix element $P_{jk}$ becomes
\begin{equation}\label{eq:6-187}
P_{jk} = \int_\eta^1 \frac{\mathscr{W}_{l+\frac{1}{2}}(\alpha_j r)\,c(r)}{\sqrt{r}}\frac{\mathscr{W}_{l+\frac{1}{2}}(\alpha_k r)}{\sqrt{r}}\frac{r^2\,dr}{r^2}
+ \int_\eta^1 r^{l+\frac{1}{2}}c(r)\frac{\mathscr{W}_{l+\frac{1}{2}}(\alpha_k r)}{\sqrt{r}}\,r^2\,dr.
\end{equation}
We can, therefore, write
\begin{equation}\label{eq:6-188}
P_{jk} = \delta_{jk}\alpha_j^2 N_{jj}\,(1-h),
\end{equation}

The characteristic equation~\eqref{eq:6-154} now becomes
\begin{equation}\label{eq:6-189}
\bigl|\delta_{jk}\alpha_j^2 N_{jj}(1-h) - l(l+1)\mathcal{C}\,Q_{jk}\bigr| = 0.
\end{equation}
In the first approximation, equation~\eqref{eq:6-189} gives
\begin{equation}\label{eq:6-190}
l(l+1)\mathcal{C} = \frac{\alpha_1^4 N_{11}(1-h)}{Q_{11}}.
\end{equation}

Comparing this formula with the corresponding formula~\eqref{eq:6-186} for the
case $b = c = 1$, we find that in this approximation
\begin{equation}\label{eq:6-191}
\mathcal{C}_{b=1} = (1-h)\mathcal{C}_{b=c=1}.
\end{equation}

It appears that in the mantle of the earth the value of gravity remains
approximately constant. For this reason the model,
\begin{equation}\label{eq:6-192}
\gamma = \text{constant}\quad\text{and}\quad c(r) = r^{-1},
\end{equation}
has been considered in some detail by Lyttkens. The factors
$(1-h)/N_{jj}$, which must be applied to the values listed in Tables~XXII--XXV to obtain the corresponding numbers for the model~\eqref{eq:6-192}, are
given in Table~XXVII. For case~(i), the derived values of $\mathcal{C}_l$ are illus\-trated in Fig.~59.

A constant value of gravity is strictly inconsistent with the model of
a mantle of constant density overlying a core; for the most general
behaviour of $g(r)$ compatible with this model is that given by equation~\eqref{eq:6-2}.
A variation of gravity compatible with this equation and which in
the interval $(0, 1)$ does not depart, sensibly, from constancy is given by
\begin{equation}\label{eq:6-193}
\gamma = \text{constant}\Bigl(\frac{r}{r_0}\Bigr),\quad(1 \leq r \leq 1).
\end{equation}

\begin{table}[ht]
\centering
\caption{The factor $(1-h)/N_{jj}$}
\label{tab:XXVII}
\begin{tabular}{c c c c c c}
\hline
$l$ & $\eta = 0.200$ & $\eta = 0.375$ & $\eta = 0.500$ & $\eta = 0.625$ & $\eta = 0.800$ \\
\hline
1 & 0.906 & 0.773 & 0.721 & 0.636 & 0.606 \\
2 & 0.930 & 0.826 & 0.738 & 0.691 & 0.626 \\
3 & 0.949 & 0.862 & 0.756 & 0.716 & 0.639 \\
4 & 0.960 & 0.884 & 0.771 & 0.731 & 0.648 \\
5 & 0.968 & 0.900 & 0.784 & 0.743 & 0.656 \\
6 & 0.973 & 0.912 & 0.795 & 0.753 & 0.662 \\
8 & 0.981 & 0.929 & 0.813 & 0.769 & 0.672 \\
10 & 0.986 & 0.941 & 0.827 & 0.782 & 0.680 \\
12 & 0.989 & 0.949 & 0.838 & 0.793 & 0.687 \\
14 & 0.991 & 0.955 & 0.847 & 0.801 & 0.693 \\
\hline
\end{tabular}
\end{table}

\figurePlaceholder{59}{The Rayleigh number $\mathcal{C}_l$ for the onset of convection in the mantle as a spherical harmonic disturbance of order $l$ for various thicknesses of the mantle. For the case illustrated here, gravity has been assumed constant in the mantle (so that $c(r)$ varies as $r^{-1}$). The different curves are distinguished by the value of the fraction ($\eta$) of the radius of the sphere occupied by the core.}

The corresponding expression for $c(r)$ is
\begin{equation}\label{eq:6-194}
c(r) = \frac{r'}{r^2},\quad (1 \leq r' \leq 1).
\end{equation}
Lyttkens has derived the characteristic numbers $\mathcal{C}_l$ for this model both
on the first and the second approximations. His results are given in
Table~XXVIII. For comparison, the results for $r = \eta$ for the model
$c(r) = r^{-1}$ are also included in this table.

\begin{table}[ht]
\centering
\caption{The characteristic numbers $\mathcal{C}_l$ for $\eta = 0.5$ and $l$ for different laws of variation
of gravity\newline
(The surfaces at $r = 1$ and $r = \eta$ both free)}
\label{tab:XXVIII}
\begin{tabular}{c c c c c c}
\hline
$l$ & \multicolumn{2}{c}{$c = r^{-1}$} & \multicolumn{2}{c}{$c = r'/r^2$} & Third \\
 & First & Second & First & Second & model \\
\hline
1 & $4.41\times 10^3$ & & $4.33\times 10^3$ & & \\
2 & $1.92\times 10^3$ & & $1.88\times 10^3$ & & \\
3 & $1.27\times 10^3$ & & $1.25\times 10^3$ & & \\
4 & $1.04\times 10^3$ & & $1.02\times 10^3$ & & \\
5 & $9.61\times 10^2$ & & $9.42\times 10^2$ & & \\
6 & $9.48\times 10^2$ & & $9.29\times 10^2$ & & \\
8 & $1.03\times 10^3$ & & $1.01\times 10^3$ & & \\
10 & $1.24\times 10^3$ & & $1.21\times 10^3$ & & \\
12 & $1.54\times 10^3$ & & $1.50\times 10^3$ & & \\
\hline
\end{tabular}
\end{table}

Another model, which is of some interest as a limiting case, is that of a
mantle overlying a core so massive that the contribution of the mantle
to the gravity can be ignored. In this case,
\begin{equation}\label{eq:6-195}
\gamma = \text{constant}\cdot r^{-2} \quad\text{and}\quad c(r) = r^{-3}.
\end{equation}
Lyttkens's results for this model are included in Table~XXVIII.

From an examination of the tabulated results it is clear that the general
nature of the dependence of $\mathcal{C}_l$ on $l$ and $\eta$ is the same for all the models.

\subsection{The case $b(r) = r^{-1}$}\label{sec:60c}
In this case, the variation of gravity is the same as in a homogeneous
sphere, but the allowed non-constancy of $b(r)$ permits a departure of the

temperature gradient from that which obtains for a uniform distribution
of heat sources in a fluid sphere. In this connexion, a case of particular
interest is when all the heat sources are confined to the core and none
exists in the mantle. Then, the general law~\eqref{eq:6-3} gives
\begin{equation}\label{eq:6-196}
\frac{d^2 T}{dr^2} + \frac{2}{r}\frac{dT}{dr} = 0.
\end{equation}
The corresponding expression for $b(r)$ is
\begin{equation}\label{eq:6-197}
b(r) = r^{-2}.
\end{equation}
Now, in virtue of equations~\eqref{eq:6-155} and \eqref{eq:6-188},
\begin{equation}\label{eq:6-198}
P_{jk} = \delta_{jk}\alpha_j^2 N_{jj}c(1-k),
\end{equation}
and the characteristic equation~\eqref{eq:6-154} becomes
\begin{equation}\label{eq:6-199}
\bigl|\delta_{jk}\alpha_j^2 N_{jj}c(1-k) - l(l+1)\mathcal{C}\,Q_{jk}\bigr| = 0.
\end{equation}
Since $Q_{jk}$ is symmetric, the secular determinant obtained from~\eqref{eq:6-199} by
transposing the rows and the columns is the same as the determinant
\eqref{eq:6-199} with $c$ replaced by $b$. Hence,
\begin{equation}\label{eq:6-200}
\text{the characteristic numbers } \mathcal{C}_l \text{ for the case~\eqref{eq:6-197}} \text{ are therefore the same
as for the case~\eqref{eq:6-192}} \text{ considered in \S\,\ref{sec:60b} above.}
\end{equation}

\section{The effect of rotation on thermal instability in a fluid sphere. The formulation of the problem}\label{sec:61}
The theory of thermal instability in a fluid sphere finds its most
important applications to the liquid core of the earth. As the core
partakes of the earth's rotation and is, moreover, the seat of the earth's
magnetic field, it is clear that the scheme of the complete problem must
allow for the simultaneous presence of rotation and magnetic field. We
have seen that an allowance for these effects in a complete manner even
in a plane geometry is a task in which the superposed temperature
gradient is radial, the imposition of symmetry of rotation generally
introduces, makes for additional analytical difficulty. There
is, indeed, no difficulty in writing down the equations of requisite
generality; the difficulty is in obtaining solutions which are valid for
the ranges of the parameters which the physical problems demand.
Only when the perturbations are axisymmetric has some advance been
made in these respects. For this reason, we shall restrict the discussion

\subsection{The representation of an axisymmetric solenoidal vector field}\label{sec:61a}
When a solenoidal vector field is axisymmetric, the restriction from
the outset to a spherical harmonic analysis, which the representation in
terms of the particular poloidal and toroidal fields described in \S\,\ref{sec:56d}
implies, can be relaxed; and this has the formal advantage.

For definiteness suppose that it represents an axisymmetric solenoidal
vector field. We shall express it as a representation of a poloidal and a
toroidal field in terms of two scalar functions $U$ and $V$ in the manner
\begin{equation}\label{eq:6-201}
\boldsymbol{u} = \operatorname{curl}\operatorname{curl}(\hat{\boldsymbol{r}}\,U) + \operatorname{curl}(\hat{\boldsymbol{r}}\,V),
\end{equation}
where $\boldsymbol{u}$, $\varpi$, and $\boldsymbol{z}$ define a system of cylindrical polar coordinates (with
the axis of symmetry in the $z$-direction), $\mathbf{1}_\varpi$, $\mathbf{1}_\varphi$, and $\mathbf{1}_z$ are unit vectors
along the three principal directions, and the defining scalars $U$ and $V$
are axially independent.

As is evident from equation~\eqref{eq:6-201}, the field derived from the poloidal
scalar $U$ has non-vanishing components only in the meridional planes;
$U$ is in fact Stokes's stream function for motions in these planes. The
field derived from the toroidal scalar $V$ has non-vanishing components
only in the transverse $\varphi$-direction. Thus $U$ defines motions which are
entirely meridional while $V$ defines motions which are entirely rotational.
When we consider the vorticity of an axisymmetric solenoidal velocity
field, we find a certain reciprocity: a poloidal velocity field gives rise to
a toroidal vorticity and conversely. Specifically,
\begin{equation}\label{eq:6-202}
\operatorname{curl}\boldsymbol{u} = -\mathbf{1}_\varphi\,\Delta_3 U + \operatorname{curl}\operatorname{curl}(\hat{\boldsymbol{r}} V),
\end{equation}
where
\begin{equation}\label{eq:6-203}
\Delta_3 = \frac{\partial^2}{\partial\varpi^2} - \frac{1}{\varpi}\frac{\partial}{\partial\varpi} + \frac{\partial^2}{\partial z^2} = \nabla^2 - \frac{1}{\varpi^2} - \frac{2}{\varpi}\frac{\partial}{\partial\varpi}
\end{equation}
is the $\varphi$-Laplacian operator for axisymmetric functions in a five-dimensional
Euclidean space.

An equivalent way of writing equations~\eqref{eq:6-201} and \eqref{eq:6-202} is
\begin{equation}\label{eq:6-204}
\boldsymbol{u} = \mathbf{1}_\varpi \times \nabla U + \operatorname{curl}(\mathbf{1}_\varphi\,\varpi V),
\end{equation}
and
\begin{equation}\label{eq:6-205}
\operatorname{curl}\boldsymbol{u} = -\mathbf{1}_\varphi\,\varpi\Delta_3 U + \operatorname{curl}(\mathbf{1}_\varpi\times\nabla V).
\end{equation}

The particular advantage of the representation~\eqref{eq:6-201} (or~\eqref{eq:6-204}) in
the present connexion is that we can write down without any difficulty

the result of the operation, any number of times, of curl on $\boldsymbol{u}$. Thus, by
repeated application of the formula~\eqref{eq:6-202}, we obtain,
\begin{equation}\label{eq:6-206}
\operatorname{curl}^{2n}\boldsymbol{u} = (-1)^n\mathbf{1}_\varphi\varpi(\Delta_3)^n U + \operatorname{curl}\operatorname{curl}[\hat{\boldsymbol{r}}(\Delta_3)^{n-1}V],
\end{equation}
\begin{equation}\label{eq:6-207}
\operatorname{curl}^{2n+1}\boldsymbol{u} = (-1)^n\mathbf{1}_\varphi\varpi(\Delta_3)^n\Delta_3 U + \operatorname{curl}[\hat{\boldsymbol{r}}(\Delta_3)^n V],
\end{equation}
and so on.

In spherical polar coordinates, the representation~\eqref{eq:6-201} takes the
form
\begin{equation}\label{eq:6-208}
\boldsymbol{u} = -\frac{1}{r}\frac{\partial^2(rU)}{\partial r\,\partial\theta}\hat{\boldsymbol{\theta}} + \frac{1}{r}\mathscr{L}^2 U\,\hat{\boldsymbol{r}} + V\hat{\boldsymbol{\varphi}},
\end{equation}
where $\mathbf{1}_r$, $\mathbf{1}_\theta$, and $\mathbf{1}_\varphi$ are unit vectors along the $r$, $\theta$, and $\varphi$
in the three principal directions. Similarly, the expression for the
Laplacian $\Delta_3$ in terms of $r$ and $\theta$ is (see equation~\eqref{eq:6-25})
\begin{equation}\label{eq:6-209}
\Delta_3 = \frac{\partial^2}{\partial r^2} + \frac{2}{r}\frac{\partial}{\partial r} - \frac{\mathscr{L}^2}{r^2} = \nabla^2 - \frac{1}{r^2\sin^2\theta}.
\end{equation}

While the nature of the functions $U$ and $V$ are not specified except
that they are axially independent, it is clear that on this theory the
basis expansions will be in terms of the fundamental solutions of the
equation
\begin{equation}\label{eq:6-210}
\Delta_3\phi = -k^2\phi,
\end{equation}
which are given by
\begin{equation}\label{eq:6-211}
\phi = j_{l+\frac{1}{2}}(kr)\frac{C_l^{3/2}(\mu)}{\sqrt{r}},
\end{equation}
where $j_{l+\frac{1}{2}}(kr)$ denotes the Bessel function of the order $l+\frac{1}{2}$, and
$C_l^{3/2}(\mu)$ is the Gegenbauer polynomial defined as the coefficient of $h^l$ in
the expansion of $(1-2\mu h+h^2)^{-3/2}$ in ascending powers of $h$:
\begin{equation}\label{eq:6-212}
(1-2\mu h+h^2)^{-3/2} = \sum_{l=0}^{\infty} h^l C_l^{3/2}(\mu).
\end{equation}
In this notation $C_l^{3/2}(\mu)$ will denote the usual Legendre polynomials.
The polynomial $C_l^{3/2}(\mu)$ is a solution of the differential equation
\begin{equation}\label{eq:6-213}
(1-\mu^2)\frac{d^2 C_l^{3/2}}{d\mu^2} - 4\mu\frac{dC_l^{3/2}}{d\mu} + l(l+2)C_l^{3/2} = 0.
\end{equation}
An alternative form of this equation is
\begin{equation}\label{eq:6-214}
\frac{d}{d\mu}\bigl[(1-\mu^2)^{3/2}C_l^{3/2}\bigr] = -(l+1)(l+2)C_{l-1}^{3/2}(1-\mu^2)^{1/2}.
\end{equation}
And, finally, we may note here that the Gegenbauer polynomials satisfy
the orthogonality relation
\begin{equation}\label{eq:6-215}
\int_{-1}^{+1} C_l^{3/2}(\mu)C_m^{3/2}(\mu)(1-\mu^2)\,d\mu = \frac{(l+1)(l+2)}{2(2l+3)}\delta_{lm}.
\end{equation}

\subsection{The perturbation equations}\label{sec:61b}
Returning to the problem of thermal instability, we consider a homo\-geneous sphere of radius $R$ rotating with an angular velocity $\Omega$ about the
$z$-axis; and we shall suppose that there is a uniform distribution of heat
sources $\epsilon$ such that the superimposed temperature gradient is given by
(cf.\ equations~\eqref{eq:6-8} and \eqref{eq:6-9})
\begin{equation}\label{eq:6-216}
\frac{\partial T}{\partial r} = -\beta r,\quad\text{where}\quad\beta = \epsilon/3\kappa = \text{constant.}
\end{equation}

In treating this problem, we shall ignore the rotational flattening of
the sphere; this is justified in our present context also since the principal
effect we are seeking is that of the Coriolis acceleration term, $2\boldsymbol{\Omega}\times\boldsymbol{u}$,
in the equation of motion; for this purpose the small departure of the
equilibrium configuration from a sphere is not essential. Under these
circumstances, the relevant perturbation equations are
\begin{equation}\label{eq:6-217}
\frac{\partial\boldsymbol{u}}{\partial t} = \alpha\nabla^2\boldsymbol{u} + 2\boldsymbol{\Omega}\times\boldsymbol{u} + \alpha\beta\gamma r\Theta\hat{\boldsymbol{r}} - \frac{1}{\rho}\nabla(\delta p),
\end{equation}
and
\begin{equation}\label{eq:6-218}
\frac{\partial\Theta}{\partial t} = \beta r u_r + \kappa\nabla^2\Theta,
\end{equation}
where $\gamma = \frac{4}{3}\pi\rho G$; and the velocity field is, of course, solenoidal.
As we have stated, we shall restrict our discussion to the case when all
the quantities describing the perturbations are axisymmetric; in par\-ticular, we shall suppose that the velocity field is specified in terms of
two scalar functions $U$ and $V$ in the manner~\eqref{eq:6-201}.

We now eliminate the term $\operatorname{grad}(\delta p/\rho)$ in equation~\eqref{eq:6-217} by taking its
curl. We then obtain:
\begin{equation}\label{eq:6-219}
\frac{\partial}{\partial t}\operatorname{curl}\boldsymbol{u} = \nu\nabla^2\operatorname{curl}\boldsymbol{u} + 2\Omega\frac{\partial\boldsymbol{u}}{\partial z} + \alpha\beta\gamma\operatorname{curl}(r\Theta\hat{\boldsymbol{r}}).
\end{equation}
It can be readily verified that
\begin{equation}\label{eq:6-220}
\operatorname{grad}(r\hat{\boldsymbol{r}}\times\boldsymbol{r}) = \mathbf{1}_\varphi\times\nabla(\varpi\Theta) + 2\Theta\operatorname{curl}(\mathbf{1}_\varphi\times\boldsymbol{r}).
\end{equation}
Also,
\begin{equation}\label{eq:6-221}
\operatorname{curl}(\hat{\boldsymbol{r}}\times\mathbf{1}_\varphi) = \mathbf{1}_\varphi \times\nabla\Bigl(\frac{1}{r}\frac{\partial}{\partial\theta}\Bigr) + \frac{1}{r\sin\theta}\operatorname{curl}\bigl(\mathbf{1}_\varphi\times\nabla\Theta\bigr).
\end{equation}

Making use of equations~\eqref{eq:6-206}, \eqref{eq:6-207}, \eqref{eq:6-220}, and \eqref{eq:6-221}, we can rewrite
equation~\eqref{eq:6-219} in the form
\begin{equation}\label{eq:6-222}
-\mathbf{1}_\varphi\varpi\frac{\partial}{\partial t}(\Delta_3 U) + \operatorname{curl}\operatorname{curl}\Bigl[\hat{\boldsymbol{r}}\frac{\partial V}{\partial t}\Bigr] = \nu\nabla^2\bigl[-\mathbf{1}_\varphi\varpi\Delta_3 U + \operatorname{curl}\operatorname{curl}(\hat{\boldsymbol{r}}V)\bigr] + 2\Omega\frac{\partial}{\partial z}\bigl[\operatorname{curl}\operatorname{curl}(\hat{\boldsymbol{r}}U) + \operatorname{curl}(\hat{\boldsymbol{r}}V)\bigr] + \alpha\beta\gamma\operatorname{curl}(r\Theta\hat{\boldsymbol{r}}).
\end{equation}
The poloidal and the toroidal fields represented in this equation must
separately vanish. Thus, we must have
\begin{equation}\label{eq:6-223}
\frac{\partial}{\partial t}\Delta_3 U = \nu\Delta_3^2 U + 2\Omega\frac{\partial V}{\partial z} + \alpha\beta\gamma\mathscr{L}^2\Theta
\end{equation}
and
\begin{equation}\label{eq:6-224}
\frac{\partial V}{\partial t} = \nu\Delta_3 V - 2\Omega\frac{\partial U}{\partial z}.
\end{equation}
Making use of the relation
\begin{equation}\label{eq:6-225}
\boldsymbol{u}\cdot\hat{\boldsymbol{r}} = -\frac{1}{r^2}\mathscr{L}^2 U,
\end{equation}
we can rewrite equation~\eqref{eq:6-217} in the form
\begin{equation}\label{eq:6-226}
\frac{\partial\Theta}{\partial t} = -\frac{\beta}{r}\mathscr{L}^2 U + \kappa\nabla^2\Theta.
\end{equation}

Seeking solutions of equations~\eqref{eq:6-223}, \eqref{eq:6-224}, and \eqref{eq:6-226} which have the
time dependence $e^{pt}$, we obtain
\begin{equation}\label{eq:6-227}
(\Delta_3 - p/\nu)\Delta_3 U = -\frac{2\Omega}{\nu}\frac{\partial V}{\partial z} - \frac{\alpha\beta\gamma}{\nu}\mathscr{L}^2\Theta,
\end{equation}
\begin{equation}\label{eq:6-228}
(\Delta_3 - p/\nu)V = 2\Omega\frac{\partial U}{\partial z},
\end{equation}
and
\begin{equation}\label{eq:6-229}
(\nabla^2 - p/\kappa)\Theta = \frac{\beta}{r}\mathscr{L}^2 U,
\end{equation}
where it should be remembered that $p$ can be complex.

By measuring distance in units of the radius $R$ of the configuration,
and making the transformations
\begin{equation}\label{eq:6-230}
\boldsymbol{v} = \frac{R^2}{\nu}\boldsymbol{u}\quad\text{and}\quad V = \frac{R^2}{\nu}\bar{V},
\end{equation}
we can express equations~\eqref{eq:6-227}--\eqref{eq:6-229} in the more convenient forms
\begin{equation}\label{eq:6-231}
(\Delta_3 - \sigma)\Delta_3 U = -\mathscr{T}\frac{\partial V}{\partial z} - \mathcal{C}\mathscr{L}^2\Theta,
\end{equation}
and
\begin{equation}\label{eq:6-232}
(\Delta_3 - \sigma)V = \mathscr{T}\frac{\partial U}{\partial z},
\end{equation}
and
\begin{equation}\label{eq:6-233}
(\nabla^2 - p\sigma)\Theta = \frac{\beta}{r}\mathscr{L}^2 U,
\end{equation}
where
\begin{equation}\label{eq:6-234}
\sigma = pR^2/\nu,\quad \mathscr{T} = 2\Omega R^2/\nu,\quad\text{and}\quad \mathcal{C} = \alpha\beta\gamma R^4/\kappa\nu.
\end{equation}

\subsection{The boundary conditions}\label{sec:61c}
Solutions of equations~\eqref{eq:6-231}--\eqref{eq:6-233} must be sought which satisfy
certain boundary conditions. These have been stated in \S\,\ref{sec:56c} (in terms
of the scalars $U$ and $V$ can be inferred from equation~\eqref{eq:6-208}) which
relates the components $u_r$, $u_\theta$, and $u_\varphi$ on the scalars $U$ and $V$.

Since $u_r$ must vanish on the boundary in all cases, we must require that
\begin{equation}\label{eq:6-235}
U = 0 \quad\text{for}\quad r = 1.
\end{equation}
If the bounding surface is rigid, then the remaining components $u_\theta$ and $u_\varphi$
must also vanish on it; and from equation~\eqref{eq:6-208} it follows that
\begin{equation}\label{eq:6-236}
\frac{\partial U}{\partial r} = 0\quad\text{and}\quad V = 0\quad\text{if the surface}\quad r = 1\quad\text{is rigid.}
\end{equation}
(Actually, the vanishing of $u_\varphi$ requires only that $(\partial V/\partial r) = 0$; but
since $V = 0$ on the boundary, the condition stated follows.) On the
other hand, if the bounding surface is free, then the vanishing of the
tangential viscous stresses on this surface requires (cf.\ equation~\eqref{eq:6-47})
\begin{equation}\label{eq:6-237}
\frac{\partial^2 U}{\partial r^2} = 0,\quad\frac{\partial}{\partial r}\Bigl(\frac{V}{r}\Bigr) = 0,
\end{equation}
\begin{equation}\label{eq:6-238}
\Theta = 0 \quad\text{for}\quad r = 1.
\end{equation}
According to equation~\eqref{eq:6-208}, these conditions are equivalent to
\begin{equation}\label{eq:6-239}
\frac{\partial}{\partial r}\Bigl[\frac{1}{r^2}\frac{\partial(r^2 U)}{\partial r}\Bigr] = \frac{\partial V}{\partial r} + \frac{2V}{r} = 0 \quad\text{for}\quad r = 1.
\end{equation}
Since $U = 0$ for $r = 1$, the first of these conditions can also be expressed
as
\begin{equation}\label{eq:6-240}
\frac{d^2 U}{dr^2}(r=1) = 0.
\end{equation}

Collecting all these conditions, we have
\begin{equation}\label{eq:6-241}
U = 0,\quad \frac{dU}{dr} = 0,\quad V = 0,\quad\text{and}\quad \Theta = 0\quad\text{for}\quad r = 1,
\end{equation}
if the bounding surface is rigid; and
\begin{equation}\label{eq:6-242}
U = 0,\quad \frac{d^2 U}{dr^2} = 0,\quad V = 0, \quad\text{and}\quad \Theta = 0\quad\text{for}\quad r = 1,
\end{equation}
if the bounding surface is free.

\subsection{The variational principle}\label{sec:61d}
The solution of equations~\eqref{eq:6-231}--\eqref{eq:6-233} together with the boundary
conditions~\eqref{eq:6-241} or \eqref{eq:6-242} constitutes a characteristic-value problem. We
shall now show that this problem can be formulated in terms of a varia\-tional principle.

Multiply equation~\eqref{eq:6-231} by $\hat{U}$ and integrate over the volume of the
unit three-dimensional sphere. Both sides of the equation can be
reduced by integration by parts and we find
\begin{equation}\label{eq:6-243}
\int\bigl[\operatorname{grad} U\cdot\operatorname{grad}\hat{U}+\sigma U\cdot\hat{U}\bigr]\,d^3x = \int U\bigl[\hat{U}\,\Delta_3 - \sigma\hat{U}\bigr]\varpi^2\,d\varpi\,dz,
\end{equation}
the integrated parts vanishing on account of the boundary conditions on
$U$. It will be observed that on the right-hand side the integration is over
the unit five-dimensional sphere.

Making use of equation~\eqref{eq:6-232}, we can rewrite~\eqref{eq:6-243} in the
manner:
\begin{equation}\label{eq:6-244}
C\int\bigl[\operatorname{grad}^2 U + \sigma|\nabla U|^2\bigr]\varpi^2\,d\varpi\,dz = \int\Bigl[\hat{U}\Bigl(\frac{\partial V}{\partial z}\Bigr)^2 + \sigma V^2\Bigr]\varpi^2\,d\varpi\,dz.
\end{equation}

Considering first the term in $V$ on the right side of this equation,
we can reduce it by making use of equation~\eqref{eq:6-232} and a sequence of
integrations by parts. Thus,
\begin{equation}\label{eq:6-245}
\int\int \hat{U}\Bigl(\frac{\partial V}{\partial z}\Bigr) \varpi^2\,d\varpi\,dz = \int\int V\Bigl(\frac{\partial\hat{U}}{\partial z}\Bigr)\varpi^2\,d\varpi\,dz,
\end{equation}

the integrated parts do not make any contributions in view of the
conditions that $U$ and either $V$ or $\partial V/\partial r$ vanish on the boundary. Con\-sidering next the term in $\Delta_3 U$ we can reduce it by the same sequence of
transformations which was used in the reduction of $\nabla_3 V$ in~\eqref{eq:6-244}. We
find
\begin{equation}\label{eq:6-246}
\int\int U_3\,\Delta_3(1-\mu^2)\,d\varpi\,dz = \int\int |\operatorname{grad} U|^2(1-\mu^2)\,d\varpi\,dz.
\end{equation}
Finally, considering the term in $\Delta U$ and letting
\begin{equation}\label{eq:6-247}
\Delta_3 U = X,
\end{equation}
we have
\begin{equation}\label{eq:6-248}
\int\int X_3\,X(1-\mu^2)\,d\varpi\,dz = \int U\,X\,\Delta_3(1-\mu^2)\,d\varpi\,dz,
\end{equation}
and the surface integral in the last line
of equation~\eqref{eq:6-247} will vanish. On the other hand, on a free boundary
\begin{equation}\label{eq:6-249}
\frac{\partial^2 U}{\partial r^2} + \frac{2}{r}\frac{\partial U}{\partial r} = 0.
\end{equation}
Hence, in all cases
\begin{equation}\label{eq:6-250}
(\mathbf{X},\mathbf{X}) = (\Delta_3 U, \Delta_3 U).
\end{equation}
and
\begin{equation}\label{eq:6-251}
\int U^2\,\Delta_3(1-\mu^2)\,d\varpi\,dz = \int |\nabla U|^2(1-\mu^2)\,d\varpi\,dz.
\end{equation}

Finally, combining equations~\eqref{eq:6-243}, \eqref{eq:6-244}, \eqref{eq:6-246}, and \eqref{eq:6-251}, we obtain
\begin{equation}\label{eq:6-252}
C\int\bigl[(\Delta_3 U)^2 + \sigma|\nabla U|^2\bigr]\varpi^2\,d\varpi\,dz = \int\bigl[(\Delta_3 U)^2 + \sigma V^2 + \mathscr{T}(\operatorname{grad} V)^2\bigr](1-\mu^2)\,d\varpi\,dz - 2\int\frac{(\Delta^2 U)^2}{U}(1-\mu^2)\,d\varpi\,dz.
\end{equation}
It can be readily shown by methods with which we are now familiar that
formula~\eqref{eq:6-252} provides the basis for a variational treatment of the
problem. In applying the variational method, it is, however, necessary
to remember that equations~\eqref{eq:6-231} and \eqref{eq:6-232} must be satisfied. Speci\-fically, one makes some assumptions concerning $U$ which are compatible
with the boundary conditions on it; but whatever assumption one makes
concerning $U$, one must determine and find $V$ as solutions of equations~\eqref{eq:6-231}
and~\eqref{eq:6-232} which satisfy the necessary boundary conditions.

\subsection{The thermodynamic significance of the variational principle}\label{sec:61e}
The physical content of the variational principle which derives from
equation~\eqref{eq:6-252} is the same as in the other cases we have considered.
However, since we have retained the time dependence through $\sigma$ and
allowed for the possibility of overstability, we must, in accordance with
the discussion in Chapter~III, \S\,23\,($b$), expect that equation~\eqref{eq:6-252}
expresses the equality
\begin{equation}\label{eq:6-253}
\epsilon_\nu = \rho\nu\!\int |\operatorname{curl}\boldsymbol{u}|^2\,dV = \epsilon_T,
\end{equation}
where $\epsilon_\nu$ and $\epsilon_T$ are the rates at which energy is, respectively, dissipated
by viscosity and liberated by the buoyancy force in the fluid sphere.

In the present instance, all the quantities occurring in~\eqref{eq:6-253} can be
evaluated directly. Thus,
\begin{equation}\label{eq:6-254}
\epsilon_\nu = \rho\nu\!\int |\boldsymbol{u}\cdot\operatorname{curl}\boldsymbol{u}|\,dV.
\end{equation}
Using the expressions~\eqref{eq:6-206} and \eqref{eq:6-208} for curl$\,\boldsymbol{u}$ and $\boldsymbol{u}$, we have
\begin{equation}\label{eq:6-255}
\epsilon_\nu = \rho\nu\int l(l+1)\bigl[(\Delta_3 U)^2\bigl]\varpi^2\,d\varpi\,dz + \rho\nu\int(\nabla V)^2\,\varpi^2\,d\varpi\,dz.
\end{equation}
After integrating by parts each of the three terms on the right-hand side,
we obtain
\begin{equation}\label{eq:6-256}
\epsilon_\nu = \rho\nu\!\int\bigl[(\Delta_3 U)^2 + |\operatorname{grad} V|^2\cdot(1-\mu^2)\bigr]\,d\varpi\,dz.
\end{equation}

\section{On the effect of rotation on the onset of stationary convection in a fluid sphere}\label{sec:62}
In this section we shall restrict our discussion to the consideration
of modes of marginal instability which are stationary (i.e.\ $\sigma = 0$).
Whether modes of overstability occur at lower values of $\mathcal{C}$ for given
$\mathscr{T}$ is an open question at the present time. However, the restriction
is not too severe: we know that when $\mathscr{T} = 0$ the onset of instability
is necessarily as a stationary mode; and we may expect that this will
continue to be the case at least for moderate values of $\mathscr{T}$.

With $\sigma = 0$, equations~\eqref{eq:6-231}--\eqref{eq:6-233} become
\begin{equation}\label{eq:6-257}
\Delta_3^2 U = -\mathscr{T}\frac{\partial V}{\partial z} - \mathcal{C}\mathscr{L}^2\Theta,
\end{equation}
\begin{equation}\label{eq:6-258}
\Delta_3 V = \mathscr{T}\frac{\partial U}{\partial z},
\end{equation}
and
\begin{equation}\label{eq:6-259}
\nabla^2\Theta = \frac{\beta}{r}\mathscr{L}^2 U.
\end{equation}
Also, the variational formula~\eqref{eq:6-252} gives
\begin{equation}\label{eq:6-260}
\mathcal{C}\!\int (\nabla\Theta)^2\,\varpi^2\,d\varpi\,dz = \int\bigl[(\Delta_3 U)^2 + \mathscr{T}^2|\nabla V|^2\bigr](1-\mu^2)\,d\varpi\,dz.
\end{equation}

\section{Remarks on geophysical applications}\label{sec:63}
We shall conclude this chapter with some brief remarks on the
geophysical applications of the theory which has been developed.

The theory of thermal convection in a sphere or a spherical shell
has an obvious bearing on problems relating to the earth's interior.
The earth's core, consisting of liquid iron, extends from about $0.55$
of the earth's radius to the centre. The mantle of the earth, which
is solid, extends from the core boundary to the surface. While the
mantle is primarily solid, it is possible that on geological time scales
it may behave as a very viscous fluid.

The theory developed in the preceding sections provides a basis for
estimating the Rayleigh number for the onset of convection and the
pattern of convection which might prevail. For the earth's core, the
most relevant quantity is the critical Rayleigh number for the onset
of convection in a sphere with a rigid boundary, modified by the
effects of rotation. The Taylor number for the earth's core is extremely
large ($\sim 10^{20}$), and it is clear that rotation will have a very
significant effect on the pattern and onset of convection.

For the mantle, the problem is that of convection in a spherical shell.
The results given in Tables~XXII--XXV show how the critical Rayleigh
number and the preferred mode of convection depend on the ratio
$\eta$ of the inner to the outer radius of the shell. For $\eta \approx 0.55$
(appropriate to the earth), the preferred mode involves spherical harmonics
of relatively low order ($l \approx 3$--$6$ depending on the boundary conditions).

The implications of these results for mantle convection, continental
drift, and the generation of the earth's magnetic field remain topics of
active research and debate.

\section*{BIBLIOGRAPHICAL NOTES}
\addcontentsline{toc}{section}{Bibliographical Notes}

\noindent\S\S\,55--59. The basic theory for the onset of thermal instability in a
fluid sphere is due to:
\begin{enumerate}
\item S.\ \textsc{Chandrasekhar}, ``The thermal instability of a fluid sphere heated
within,'' \emph{Phil.\ Mag.}\ (7), \textbf{43}, 1317--29 (1952).
\end{enumerate}

\noindent The theory of \S\S\,55--59 follows this paper.

\noindent See also:
\begin{enumerate}\setcounter{enumi}{1}
\item S.\ \textsc{Chandrasekhar}, ``The onset of convection by thermal instability in
spherical shells,'' \emph{Phil.\ Mag.}\ (7), \textbf{44}, 233--41 (1953).
\item G.\ E.\ \textsc{Backus}, ``Critique of `The thermal instability of a fluid sphere
heated within,' '' 1953 (unpublished).
\end{enumerate}

\noindent The `exact' values listed in Tables~XX and XXI are from reference 3.

\noindent\S\,60. The treatment follows:
\begin{enumerate}\setcounter{enumi}{3}
\item S.\ \textsc{Chandrasekhar}, ``The onset of convection by thermal instability in
spherical shells,'' \emph{Phil.\ Mag.}\ (7), \textbf{44}, 233--41 (1953).
\end{enumerate}

\noindent The extension to the case of non-uniform gravity and non-uniform
distribution of heat sources is due to:
\begin{enumerate}\setcounter{enumi}{4}
\item E.\ \textsc{Lyttkens}, ``On the onset of convection by thermal instability in
spherical shells,'' \emph{Astrophys.\ J.}, in press.
\end{enumerate}

\noindent Tables~XXVII and XXVIII are from reference 5.

\noindent\S\S\,61 and 62. The treatment of the problem allowing for the effects
of rotation follows:
\begin{enumerate}\setcounter{enumi}{5}
\item S.\ \textsc{Chandrasekhar}, ``The instability of a layer of fluid heated below
and subject to the simultaneous action of a magnetic field and rotation,''
\emph{Proc.\ Roy.\ Soc.}\ (London) A, \textbf{225}, 173--84 (1954).
\item S.\ \textsc{Chandrasekhar} and D.\ \textsc{Elbert}, ``The instability of a layer of
fluid heated below and subject to Coriolis forces,'' \emph{Proc.\ Roy.\ Soc.}\
(London) A, \textbf{231}, 198--210 (1955).
\end{enumerate}

\noindent The discussion in \S\,61 on the representation of axisymmetric solenoidal vector fields
is based on:
\begin{enumerate}\setcounter{enumi}{7}
\item S.\ \textsc{Chandrasekhar}, ``On the expansion of functions which satisfy
four boundary conditions,'' \emph{Proc.\ Nat.\ Acad.\ Sci.}, \textbf{43}, 521--7
(1957).
\end{enumerate}

\noindent\S\,63. For geophysical applications see:
\begin{enumerate}\setcounter{enumi}{8}
\item W.\ M.\ \textsc{Elsasser}, ``Hydromagnetism,'' \emph{Am.\ J.\ Phys.}, \textbf{23}, 590--609 (1955);
\textbf{24}, 85--110 (1956).
\item E.\ C.\ \textsc{Bullard}, ``The stability of a homopolar dynamo,'' \emph{Proc.\
Camb.\ Phil.\ Soc.}, \textbf{51}, 744--60 (1955).
\end{enumerate}
